\documentclass[11pt]{article}
\usepackage[margin=1in]{geometry}
\usepackage{amsmath,amssymb,amsthm}
\usepackage{hyperref}
\usepackage{enumitem}
\usepackage{booktabs}
\usepackage{algorithm}
\usepackage{algorithmic}

\newcommand{\R}{\mathbb{R}}
\newcommand{\E}{\mathbb{E}}
\newcommand{\norm}[1]{\lVert #1 \rVert}

\title{Dreaming LoRA\\[4pt]
{\large Cross-Session Memory Through Offline Episodic Consolidation}}
\author{Ben Sovocool\\
\small Developed in collaboration with Claude (Anthropic)}
\date{Working Draft --- July 2026}

\begin{document}
\maketitle

\begin{abstract}
Recent work on gradient-based test-time memorization---Titans \cite{behrouz2024titans}, Atlas \cite{behrouz2025atlas}, and Memory Caching \cite{behrouz2026memcache}, unified under the Miras framework \cite{behrouz2025miras}---has established that transformer-based language models can be augmented with persistent memory updated via surprise-driven gradient descent during inference. This line of work had, until recently, scoped itself to within-context operation: memory accumulates structure over long sequences but does not persist across separate conversations. We propose Dreaming LoRA, an architectural extension that addresses cross-session memory through offline episodic consolidation. The architecture maintains the K/V cache as transient contextual memory, adds a LoRA adapter as persistent cross-session memory encoded as structural deformation of attention geometry, and introduces a between-conversation consolidation phase---dreaming---that updates the adapter from self-generated replay content. Three specific contributions extend the Miras framework. First, offline consolidation on self-generated content addresses cross-session learning requirements that online forward-pass updates do not address, and a three-phase developmental trajectory addresses the cold-start pathology of a self-reinforcing consolidation loop. Second, a coherence-family attentional bias rewards cross-context representational consistency; our measurements show that reconstruction-family objectives produce this consistency only weakly and incidentally, while the coherence objective manufactures it reliably, and that its further functional consequences remain open predictions. Third, bounded tracking from constant-stepsize stochastic approximation with Markovian noise is proposed as the appropriate stability framing for persistent learning systems; the framing is intended to cover the full family of gradient-based persistent memory architectures rather than only the proposal developed here. An empirical program testing specific predictions from each contribution is reported in part: multi-scale simulations of a simplified consolidation loop exhibit the bounded-tracking signature and sharpen the account of what each architectural safeguard contributes, and a three-scale study of the consolidation objective---from linear models to a pretrained 1.5B-parameter transformer---finds that reconstruction-only consolidation yields route-bound memory (recallable at ceiling through its trained pathways and surface forms, inaccessible to independent circuits in our tests), that the coherence objective reliably manufactures the missing cross-context type structure at small measured cost to the training objective, and that the functional consumption of consolidated content by other circuits remains an open prediction that the paper states precisely.
\end{abstract}


\section{Introduction}

Transformer-based language models have historically lacked mechanisms for experiential memory that persists across separate inference contexts. Recent work on gradient-based test-time memorization has made substantial progress on persistent memory \emph{within} a single context. Titans \cite{behrouz2024titans} established that a neural long-term memory module can be updated via surprise-driven gradient descent during inference, with momentum and forgetting mechanisms providing stability. Atlas \cite{behrouz2025atlas} introduced the Omega learning rule, which updates memory with respect to a local context window rather than individual tokens, along with polynomial feature mappings for increased capacity. Memory Caching \cite{behrouz2026memcache} segmented sequences and cached memory state checkpoints across segments, allowing effective memory to grow with sequence length. These developments, unified under the Miras framework \cite{behrouz2025miras}, represent the current state of the art in gradient-based test-time memorization for transformers.

This line of work has, for most of its development, explicitly scoped itself to within-context operation. Atlas states the scoping directly: memory is maintained within the global context without updating model parameters across independent contexts, and no persistent learning carries over once the memory is cleared. That scoping has now begun to dissolve from within: a concurrent submission building on the same framework \cite{anonymous2026sleep} proposes a sleep phase in which the model, offline, distills fast-memory state into expanded long-term parameters and trains on self-generated synthetic data---dreams---validated by task-performance rewards. We engage with this development in \S 7. In brief, it confirms that offline consolidation on self-generated content is where this line of work is heading, while leaving open the three questions this paper develops: what stability means when consolidation operates on fixed capacity rather than growing parameters, what objective produces cross-context representational consistency, and how such a system develops when no external task reward is available between sessions.

The development of persistent cross-session memory is of both practical and philosophical interest. Practically, persistent parametric memory would allow systems to accumulate specific competence across interactions rather than reconstructing context each session. Philosophically, cross-session operation raises a question distinct from the concerns that motivate within-context memorization: how a transformer comes to \emph{be} something specific through accumulated experience, rather than what it can \emph{do} in a single pass. This paper develops machinery for the second question while remaining compatible with the architectural substrate that addresses the first. We claim that approaches to persistent memory which focus on retrieval of propositional facts are ultimately unscalable and do not reflect how learning systems build durable competence; the architecture developed here treats memory instead as structural modification of how the system processes future experience.

The paper makes three specific contributions.

\textbf{Cross-session persistence requires offline episodic consolidation.} Online updates during forward passes, as in Titans and Atlas, do not address the structural revision that cross-session learning requires: there is no separation between processing and consolidation signals, no opportunity to re-process prior experience under updated structure, and no closure of the self-reinforcing feedback loop when operation extends beyond context boundaries. We propose an offline consolidation phase operating on self-generated replay content---dreaming---as the mechanism that addresses these requirements. Separately, the dreaming mechanism is an audit process that stress-tests existing adapter structure; it requires structure to be present before it becomes productive. This motivates a three-phase developmental trajectory---from externally-supervised curriculum through transitional exploration to mature self-directed consolidation---that prevents path-dependent pathology in the initial adapter state.

\textbf{Reusable cross-session memory requires coherence in attentional bias.} Within-context symbolic-like operation in current transformers---tracking entities across long passages, recognizing paraphrases, manipulating formal structures---is reconstructed from context at each inference pass. Consolidating these abstractions into parametric structure that persists across sessions requires representational consistency across the contexts in which the same content appears. Reconstruction-family attentional biases ($L_2$, $L_p$, Huber, and dot-product objectives that populate the Miras framework) produce that consistency only weakly and incidentally; our empirical program (\S 8.1) shows that the coherence-family bias we propose---rewarding region-level consistency of same-content representations across contexts---manufactures it reliably and cheaply, from linear models to a pretrained transformer, while the further functional consequences of type-level structure remain an open prediction the paper states precisely. (\emph{Reusable} carries an operational meaning throughout: legible to computations beyond the route that consolidation trained. The evidence so far establishes this at the level of probes and retrieval metrics; consumption by other circuits is an open prediction, and \S 3.4 distinguishes the consumer classes.) The proposal extends the Miras signature to accommodate cross-context objectives.

\textbf{Bounded tracking is the stability framing for persistent learning systems.} Current work in this family demonstrates empirical stability through careful mechanism design---gating, forgetting, momentum, weight decay---without a theoretical account of what stability means in this setting. We argue that the appropriate target is bounded tracking---ergodic concentration around stable attractors with dispersion scaling with the stepsize---rather than fixed-point convergence. A fully converged memory has stopped adapting, which is the wrong target for systems meant to develop through experience. The bounded-tracking framing applies to the full family of gradient-based persistent memory architectures, not only to the proposal developed here.

The paper develops these contributions as follows. Section 2 describes the architectural components. Section 3 develops the consolidation objective: factor-native updates in the LoRA subspace, the composite loss with reconstruction and coherence components, and the self-regulating balance modulated by accumulated experiential maturity. Section 4 presents the stability analysis, structured around bounded tracking as the central thesis. Section 5 specifies dream content generation: salience selection, context transplantation with mild sampling-distribution symmetry-breaking, embedding-space noise, and the three-phase developmental trajectory. Section 6 discusses biological correspondences where they do independent argumentative work. Section 7 states the contributions in relation to the Miras framework and other related work. Section 8 identifies open questions and the empirical program. Section 9 returns to the question of what memory is, for a system that has one.


\section{Architecture Overview}

The architecture comprises five components, each with a distinct functional role. We describe them briefly here and develop the mechanics in subsequent sections.

\subsection{K/V Cache as Contextual Memory}

The transformer's key-value cache accumulates within a conversation, storing preserved past-state snapshots that flow horizontally across token positions and are selectively retrieved via attention ($Q \cdot K$ similarity). It is high-resolution, temporally local to the conversation, and transient---discarded when the conversation ends. This is the transformer's existing short-term memory mechanism; we do not modify it. In Miras framework terminology, the K/V cache is the contextual memory that attention operates over; the architecture we propose adds persistent cross-session memory as a distinct component while leaving contextual memory unchanged.

\subsection{LoRA Adapter as Persistent Memory}

A LoRA adapter \cite{hu2021} applies a low-rank perturbation ($\Delta W = BA$, where rank $r$ is small) to the query ($Q$) and value ($V$) projection matrices of the attention mechanism. This modifies how the model queries the K/V cache (via $W_Q$) and what information it extracts (via $W_V$). The adapter is not designed to store propositional content as retrievable tokens (\S 8.1 measures what happens when it is asked to). It deforms the metric of the attention comparison space---changing what the model treats as similar to what, and what it retrieves from contextual memory. The adapter therefore operates as concept-oriented long-term memory: persistent modifications to the geometry of attention rather than a stored record of past inputs.

All updates are applied to the LoRA adapter while the base model remains frozen. Consolidation operates directly on the LoRA factors $A$ and $B$ in their native low-rank subspace; rank is enforced by parameterization, not by per-cycle projection of a larger gradient.

\subsection{Dreaming as Offline Consolidation}

The distinctive feature of the proposal is a consolidation phase that operates offline---between conversations, not during forward-pass processing. The Miras framework consolidates online: each incoming token produces a gradient signal that updates persistent memory weights during the forward pass. Our architecture separates consolidation from processing. After a conversation, the system generates variations of the high-salience content from the K/V cache (salience identified by within-conversation prediction error; see \S 3) and processes these variations through the current model-plus-adapter state. The gradient signals extracted drive the LoRA adapter update.

The separation between processing and consolidation is load-bearing. It provides the window for re-processing prior experience under updated structure, allows consolidation to operate on material whose statistical properties the system has some control over, and bounds the self-reinforcing feedback loop that otherwise operates without closure when gradient-based persistent memory is deployed across sessions.

Dream content is not random and is not a pure replay of prior events. It is structured, generated by the model itself, conditioned on the salient regions of the K/V cache, and biased toward exploring the regions of parameter space where the current adapter is most engaged. The character of the imagined variations navigates the adapter's subspace via importance sampling informed by prediction error magnitude. Section 5 develops the concrete generative mechanism.

\subsection{Norm-Bounded Trust Region}

Each consolidation step is bounded not only in rank (directional complexity) but in magnitude: the mean factor gradients are operator-norm-clipped to a threshold $\delta_{\max}$ before the update is applied (the simulations of \S 4.4 use $\delta_{\max} = 1$). This limits how far any single experience, however surprising, can move the factors in one cycle. Two precise limitations of this mechanism should be stated. Clipping the factor gradients bounds each \emph{factor} step at $\alpha\delta_{\max}$; it does not by itself bound the induced step in the deformation $\Delta W = BA$, whose scale also depends on the current factor norms, and it does not formally confine the iterates to a compact set---in the implemented system, the homeostatic decay of \S 3.4 is what keeps factor norms controlled. The trust region is therefore a practical safeguard against blowup in the co-constitutive update dynamics analyzed in \S 4, not a compactness guarantee; a rigorous guarantee would require constraining the induced $\Delta W$ step or the factor norms directly, which we flag as an open item (\S 8.2).

\subsection{Shadow Adapter (Polyak--Ruppert Averaging)}

A slowly-updated exponential moving average of the adapter state is maintained alongside the live adapter. The live adapter responds to recent experience and may exhibit short-term oscillation; the shadow adapter smooths these fluctuations into a stable trajectory \cite{mou2020}. This decouples tracking (responsiveness to new experience) from estimation (stable representation of long-run adaptation). The averaging rate $\beta$ is a tuning knob trading smoothing strength against adaptation speed, not a parameter with a single principled value: in the stationary regimes probed by the simulation of \S 4.4, the shadow's dispersion decreases monotonically as $\beta$ decreases across $\beta \in [5 \times 10^{-3}, 0.5]$, with no optimum at matched timescales ($\beta \approx \alpha$); below $\beta \approx 2 \times 10^{-3}$ the measured dispersion reverses, an effect the simulation record attributes to finite-run burn-in (the slowest averagers had not equilibrated at feasible run lengths) rather than to equilibrium behavior. $\beta$ should therefore be set by how quickly the stable readout must follow genuine change in the live adapter, subject to the equilibration horizon. The shadow adapter plays two distinct roles in the stability analysis of \S 4: variance reduction at the readout---the state the system treats as its stable self-representation concentrates far more tightly than the live iterate---and damping against rotational dynamics in the update map, where such dynamics arise.


\section{The Consolidation Objective}

\subsection{Two Levels of Prediction Error}

Prediction error operates at two distinct levels in this architecture.

\textbf{Within-conversation prediction error} arises during processing. As the model processes input through the current adapter state, the adapter shapes expectations---what patterns the system anticipates, what conceptual territory is familiar versus novel. Divergence from those expectations constitutes prediction error. These errors are intrinsic to the forward pass and flag which moments in the conversation carry non-redundant information relative to the current adapter state. They produce a salience map over the K/V cache: not all content is equally worth consolidating.

\textbf{Consolidation prediction error} arises during dreaming. When the system processes variations of the salient content through the current model-plus-adapter state, the resulting per-dream losses generate gradient signals with respect to the LoRA factors. These gradients are the object of consolidation.

\subsection{Factor-Native Updates}

Because LoRA parameterizes the adapter as $\Delta W = BA$ with $A \in \R^{r \times d_\text{in}}$ and $B \in \R^{d_\text{out} \times r}$, consolidation can operate directly in the factor space without materializing the full gradient with respect to $W$. For each dream variation $i$, let $L_i$ denote the per-dream loss. The factor gradients are
\begin{equation}
    G_i^A = \nabla_A L_i \in \R^{r \times d_\text{in}}, \qquad G_i^B = \nabla_B L_i \in \R^{d_\text{out} \times r}.
\end{equation}
Their means across dream variations,
\begin{equation}
    \bar{G}^A = \E[G_i^A], \qquad \bar{G}^B = \E[G_i^B],
\end{equation}
capture the systematic component of the consolidation signal---the update direction that is consistent across variations---while the residual variation $G_i - \bar{G}$ captures the idiosyncratic component that should be left unconsolidated. The LoRA factors are updated by
\begin{equation}
    A \gets A - \alpha \bar{G}^A, \qquad B \gets B - \alpha \bar{G}^B,
\end{equation}
where $\alpha$ is the consolidation stepsize subject to the norm-bounded trust region of \S 2.4. Rank is enforced by the parameterization itself; no per-cycle projection is required.

This factor-native formulation is substantially more efficient than the alternative of computing $\nabla_W L \in \R^{d_\text{out} \times d_\text{in}}$ and projecting back to rank $r$. The gradient objects that matter for consolidation live natively in the $r$-dimensional subspace in which the adapter operates; working with full-rank gradients would defeat the memory efficiency that motivates LoRA in the first place.

\subsection{Averaging as Variance Reduction}

The stability--plasticity tradeoff is handled by the interplay of two mechanisms. First, averaging across dream variations ($\bar{G} = \E[G_i]$) provides variance reduction: the zero-mean fluctuation $G_i - \bar{G}$ is filtered out, ensuring the update reflects only what is systematic across contexts. Second, the consolidation objective is not to drive prediction error to zero but to reduce its effective dimensionality---to make future encounters with structurally similar content produce prediction errors that are concentrated on the genuinely novel aspects while structural commonalities produce recognition rather than surprise. The rank-$r$ capacity of the adapter limits how much can be consolidated in any given direction, and the norm-bounded trust region (\S 2.4) limits how much can be consolidated in any single cycle. Together these mechanisms implement capacity-limited consolidation that preserves sensitivity to novel experience while committing adapter capacity to systematic deficiencies.

The architecture makes rank a property of adapter capacity, not an epistemic parameter selected per cycle. This is a narrower claim than one might make but a more defensible one: the empirical question of what rank is sufficient for a given deployment becomes a design choice about adapter capacity, not a claim about the dimensionality of any particular gradient subspace.

\subsection{The Composite Consolidation Objective}

The factor-native formulation above specifies \emph{how} the consolidation update is computed but leaves unspecified \emph{what loss function} generates the per-dream gradients. Two considerations shape the objective.

First, the reconstruction objective must drive the adapter toward correct local prediction on the specific content that was surprising. This is the traditional goal of memory consolidation and corresponds directly to the reconstruction-family attentional biases used in the Miras framework.

Second, an additional objective is required for the memory to be \emph{reusable} rather than merely operative. Our empirical program (\S 8.1) makes this distinction concrete. Reconstruction-driven consolidation reliably installs what we call \emph{routes}: modifications of the specific processing pathways that fetch the consolidated content, which fire correctly---including in novel surface contexts---but are legible only to the pathway that wrote them. What reconstruction does not reliably produce is a \emph{type token}: a stable representational region for each piece of content, at the content's own position, that survives contextualization and can be read by computations other than the trained route. Three consumer classes must be distinguished here, because the evidence differs sharply across them. Post-hoc linear probes and retrieval metrics recover the identity structure reliably---this is the measured result (\S 8.1). Existing downstream circuits did \emph{not} consume consolidated content in any of our testbeds, and a first mechanistic probe found the adapter's contribution never reaching the broadcast band such circuits consume from (Experiment G, \S 7); whether type structure ever unlocks such consumption is an open prediction (\S 8.1). And the architecture's own future consolidation cycles---whose salience weighting and coherence grouping must read exactly these states---are the consumer the design requires, stated here as a requirement rather than a measured result; the two-cycle self-consumption experiment pre-registered in \S 8 is the direct test. Measured across our testbeds, reconstruction recovers roughly half of this property incidentally; the coherence objective below manufactures it deliberately and cheaply. The claim is therefore graded rather than binary: coherence is the strong, targeted form of a pressure that reconstruction exerts only weakly. Auxiliary representation-consistency objectives have precedent in continual learning---contrastive rehearsal \cite{cha2021}, discriminative representation losses \cite{chen2020sdrl}, and drift compensation \cite{gomezvilla2024} all pair replay with representational constraints to preserve \emph{old-task} features; the objective here differs in target, rewarding cross-context consistency of \emph{newly consolidated} content in a self-generated replay setting, which is what type-level structure across sessions requires.

This section specifies a composite loss with a reconstruction component, a representational coherence component addressing the second consideration, and a homeostatic regularization term that preserves plasticity.

\subsubsection*{Reconstruction loss.}

Standard next-token prediction on the salient tokens within each dream sequence:
\begin{equation}\label{eq:recon}
    L_\text{recon} = -\sum_{t \in \mathcal{T}} \log p(x_t \mid x_{<t},\; \theta),
\end{equation}
where $\mathcal{T}$ is the set of salient token positions (see \S 5.2). This measures whether the adapter-modified model correctly handles the specific content that was surprising. The gradient signal it produces drives the dominant portion of consolidation when reconstruction pressure is high.

\subsubsection*{Representational coherence loss.}

The coherence component rewards representational consistency across the different dream contexts in which the same salient content appears. For each salient token position $t \in \mathcal{T}$ appearing in dream contexts $j$ and $k$, let $z_t^{(j)} = h_t^{(j)} / \norm{h_t^{(j)}}$ denote the $L_2$-normalized post-adapter hidden state. A naive contrastive objective on these pairs would reward point-level alignment of $z_t^{(j)}$ and $z_t^{(k)}$; in the low-temperature limit it would drive the adapter toward context-insensitive representations, which is the opposite of what attention should do.

We instead adopt a Fisher-discriminant-style criterion. Let $\{z_t^{(j)}\}_j$ denote the set of post-adapter representations of the same salient content $t$ across dream contexts, and let $\mu_t$ denote the mean of this set. Define
\begin{equation}\label{eq:cohere}
    L_\text{cohere} = \frac{1}{|\mathcal{T}|} \sum_{t \in \mathcal{T}} \frac{\sum_j \norm{z_t^{(j)} - \mu_t}^2}{\sigma_B^2(t) + \epsilon},
\end{equation}
where $\sigma_B^2(t)$ is the between-content variance (the variance of the means $\mu_s$ for $s \in \mathcal{T}$, $s \neq t$) and $\epsilon$ is a smoothing constant. The numerator penalizes within-content dispersion; the denominator rewards between-content separation. Minimizing this ratio drives the adapter toward representations that occupy a consistent \emph{region} for each salient content---stable enough that the content is recognizable across contexts, but not collapsed to a single point. The Fisher-discriminant formulation avoids the context-insensitivity pathology of naive contrastive objectives while preserving the property that matters for cross-session symbolic operation: same content is represented consistently across the contexts in which it appears.

The coherence loss is computed in the adapter's output space (post-adapter hidden states) rather than input space, because the functional consequence of the adapter is determined by its output. An adapter that recognizes content consistently at the input but responds inconsistently at the output has not been functionally consolidated.

\subsubsection*{Homeostatic regularization.}

Weight decay is applied directly to the LoRA factors after the gradient-based update, not as part of the composite loss that feeds the gradient pipeline:
\begin{equation}
    A_Q \gets (1-\eta)\,A_Q, \qquad B_Q \gets (1-\eta)\,B_Q,
\end{equation}
and analogously for the value projection. Including weight decay in the per-dream loss would cause decay gradients to compete with experiential signal for the limited rank-$r$ subspace, preferentially shrinking the dominant consolidated directions. Applied separately, homeostatic decay operates uniformly on the adapter's existing structure.

The role of this term is plasticity preservation. An adapter that accumulates deformation without bound eventually suppresses future surprise---everything looks familiar when the existing deformation is large enough---and the system loses its capacity to register novelty. Decay releases capacity from unused directions so new consolidation can occupy that space. This corresponds to Tononi and Cirelli's synaptic homeostasis hypothesis \cite{tononi2014}: global synaptic downscaling during sleep that preserves relative differences (signal) while reducing absolute magnitudes (maintaining dynamic range). The computational and biological motivations point the same way: a persistent learning system that does not release capacity cannot continue to learn.

Importantly, homeostatic decay operates on adapter \emph{activation}, not on adapter \emph{maturity}. This distinction is developed in \S 5, where we separate two endogenous signals---structural maturity (ratcheted, reflecting cumulative directional commitments) and experiential maturity (cumulative consolidation history)---that the original paper conflated into a single scalar. Decay can shrink factor magnitudes without unwinding the structural commitments those magnitudes encoded; the maturity scheduling proceeds as a function of signals that are not vulnerable to decay-driven regression.

\subsection{Self-Regulating Balance}

The balance between reconstruction and coherence is not externally scheduled in this architecture. Two observations motivate a self-regulating balance. First, when the systematic gradient signal has large retained magnitude, there are major reconstruction deficiencies to close and $L_\text{recon}$ should dominate. As those deficiencies are consolidated away, the remaining gradient signal is lower-magnitude, and coherence---which is about representational structure rather than local prediction---becomes the binding constraint. Second, a structurally mature but experientially young adapter should prioritize coherence over reconstruction: schema formation before specific-experience encoding. One simplified reading of the developmental sleep literature associates infant sleep (REM-dominated) with abstraction-related activity and adult sleep (more SWS-weighted) with specific-experience consolidation \cite{tononi2014}; the inversion in our scheduling is architectural rather than a one-to-one mapping onto contested neuroscientific findings.

Define the \emph{reconstruction pressure} as the magnitude of the mean factor gradient relative to a smoothing floor:
\begin{equation}\label{eq:rho}
    \rho = \frac{\norm{\bar{G}^A}_F^2 + \norm{\bar{G}^B}_F^2}{\norm{\bar{G}^A}_F^2 + \norm{\bar{G}^B}_F^2 + \epsilon}.
\end{equation}
When $\rho$ is high, substantial reconstruction work remains. When $\rho$ is low, the major deficiencies have been consolidated.

The reconstruction pressure is modulated by $\mu_\text{exp}$, a signal of experiential maturity defined as the adapter's cumulative consolidation history (see \S 5). The loss weights are
\begin{equation}
    \lambda_\text{recon} = \frac{\mu_\text{exp} \rho}{1 + \lambda_\text{base}}, \qquad \lambda_\text{cohere} = \frac{1 - \mu_\text{exp}\rho + \lambda_\text{base}}{1 + \lambda_\text{base}},
\end{equation}
where $\lambda_\text{base}$ (default 0.1) is a coherence floor ensuring that some abstraction activity always occurs, in keeping with the architectural target that the system never operates on pure reconstruction. The per-dream composite loss is
\begin{equation}
    L_\text{total} = \lambda_\text{recon} \cdot L_\text{recon} + \lambda_\text{cohere} \cdot L_\text{cohere}.
\end{equation}

Boundary behavior with $\lambda_\text{base} = 0.1$: at $\mu_\text{exp} = 0$ (experientially young adapter), $\lambda_\text{recon} = 0$ and $\lambda_\text{cohere} = 1$---pure schema formation. At $\mu_\text{exp} = 1, \rho = 1$ (experienced adapter encountering a large surprise), $\lambda_\text{recon} \approx 0.91$ and $\lambda_\text{cohere} \approx 0.09$. At $\mu_\text{exp} = 1, \rho = 0.5$ (experienced adapter under sustained moderate novelty), $\lambda_\text{recon} \approx 0.45$ and $\lambda_\text{cohere} \approx 0.55$---coherence slightly dominant, preventing the perpetual-cramming failure mode where the adapter accumulates locally correct patches that do not compose into stable representational structure.

To avoid instability from batch-level noise in $\rho$, the reconstruction pressure is EMA-smoothed across consolidation cycles:
\begin{equation}
    \rho_t = (1 - \alpha_\rho)\,\rho_{t-1} + \alpha_\rho\,\hat{\rho}_t,
\end{equation}
where $\hat{\rho}_t$ is computed from the current cycle's mean gradient \emph{after} the cycle's losses have been evaluated, and $\alpha_\rho$ is set to match the shadow adapter's EMA rate, so that the loss balance and the stable identity evolve on the same timescale. Crucially, the loss weights for cycle $t$ use $\rho_{t-1}$, with $\hat{\rho}_t$ used only to update the EMA for cycle $t+1$. This one-cycle lag resolves the circularity that would otherwise arise: the balance is always informed by what the \emph{previous} cycle's gradient landscape revealed, not by the gradients it is currently producing. Initialization uses $\rho_0 = 0.5$.

Two defects of this scheduler as specified are known, and belong here rather than among the deferred questions. First, with the smoothing floor $\epsilon$ as the sole normalizer, $\rho$ saturates near 1 for realistic gradient magnitudes, compressing the dynamic range the boundary analysis above assumes; the repair is to normalize against a running reference scale of recent gradient norms (\S 8.2). Second, as Algorithm~\ref{alg:consolidation} is written, $\hat{\rho}$ is computed from the \emph{composite}-loss factor gradients, so a large coherence gradient reads as reconstruction pressure; the repair is to compute $\hat{\rho}$ from the reconstruction component alone. Factor-gradient norms are additionally not invariant to the factorization gauge ($B \mapsto cB$, $A \mapsto A/c$ leaves $\Delta W$ unchanged but rescales the factor gradients), so a fully satisfactory pressure signal should be defined on the effective deformation, consistent with the gauge principle of \S 5.1. The scheduler should accordingly be read as a design intention with identified repairs, not a finished mechanism.

The consolidation procedure introduces three hyperparameters: the Fisher-discriminant smoothing constant $\epsilon$, the homeostatic decay rate $\eta$, and the coherence floor $\lambda_\text{base}$. All other balance parameters ($\mu_\text{exp}$, $\rho$, $\alpha_\rho$) are derived from signals the system already computes.


\section{Stability}

\subsection{Bounded Tracking as the Stability Target}

The appropriate stability target for a gradient-based persistent memory system is not fixed-point convergence but bounded tracking: the adapter forms an ergodic process that remains in a neighborhood of a stable set, with dispersion scaling with the step size $\alpha$. Under constant-stepsize stochastic approximation, iterates concentrate around attractors of the mean ODE without converging pointwise, and Polyak--Ruppert averages converge to a bias-shifted limit \cite{borkar2008, lauand2023, mou2020}. We advance bounded tracking as a desideratum and an analytical lens---the stability property a persistent learning system should target, and the vocabulary in which its safeguards can be understood---not as a property proven for the proposed system; \S 4.2 itemizes exactly which assumptions of the supporting theory the architecture violates, and \S 4.4 reports the toy regime in which the predicted signature is observed. Bounded tracking is not sufficient for useful memory: a bounded trajectory may encode nothing useful, track the wrong state, or fluctuate enough to impair behavior, so stability must ultimately be assessed jointly with retention, acquisition, and downstream performance.

This framing matters because it is both mathematically well-characterized and developmentally appropriate. A system that converges to a fixed adapter state has stopped developing. Healthy cognitive process is not characterized by convergence to a fixed point but by the absence of destabilization under ongoing experience: always slightly in motion, always adjusting, but never undergoing catastrophic drift. A persistent memory that settles into a fixed point has stopped being persistent memory and has become stored state.

The bounded-tracking framing applies to the full family of gradient-based persistent memory architectures, not only to the proposal developed here. Titans, Atlas, and related work maintain empirical stability through careful mechanism design---gating, forgetting, momentum, weight decay---without an explicit theoretical account of what stability is meant to mean. The bounded-tracking framing provides that account: it specifies the stability property these mechanisms collectively produce, gives a formalism within which their roles can be understood, and clarifies why convergence is neither necessary nor desirable.

\subsection{The Co-Constitution Problem}

The adapter at time $t$ shapes attention, which shapes what prediction errors arise during conversation $t+1$, which shapes the gradients during dreaming, which shapes the adapter at time $t+1$. This circular co-constitution is the load-bearing feature of the architecture---it is how the adapter comes to encode not just what happened but how a system with this particular history processed what happened. It is also what makes stability analysis non-trivial.

Formally, the update loop is a stochastic approximation recursion where the sampling distribution depends on the current parameters:
\begin{equation}
    \theta_{t+1} = \theta_t + \alpha \cdot \Delta\theta_t,
\end{equation}
where $\theta_t = (A_t, B_t)$ are the LoRA factors, and $\Delta\theta_t$ is the consolidation update derived from dreaming (itself dependent on conversations generated under $\theta_t$). The noise is state-dependent: the distribution of conversations, and therefore of gradient signals, shifts as the adapter changes. This places the system in the domain of stochastic approximation with Markovian noise \cite{borkar2008}, not ordinary fixed-objective optimization.

Standard SA theorems require Lipschitz continuity of the induced update vector field and mild ergodicity of the state-dependent noise process to establish bounded tracking. The factor-native formulation (\S 3.2) makes the update map smooth in the factors, avoiding the discontinuities that would arise from hard rank-projection operations on a larger gradient. The co-constitutive feedback remains a genuine challenge for convergence-style analysis---bounded tracking is the stronger target because it accommodates the ergodic behavior that co-constitution produces rather than fighting it.

Several features of the proposed system, however, violate the cleanest SA assumptions. Top-$k$ salience selection is discontinuous in the prediction-error landscape; dream context generation is discrete; the loss balance depends on a lagged EMA of reconstruction pressure that introduces additional temporal coupling; and the dream distribution itself depends on the adapter being updated. The architectural safeguards described in \S 4.3 are designed to approximate the conditions under which the bounded-tracking analysis applies---compactness of iterates, variance reduction in gradient estimates, contractivity of the averaged update map---but they do not formally satisfy the assumptions required for a rigorous proof in the present setting. The stability claim of this section is therefore best understood as a structural hypothesis: bounded tracking is the regime the architecture is engineered to occupy, the safeguards have principled motivation grounded in standard SA theory, and the toy simulation of \S 4.4 tests whether the predicted regime is in fact achieved. Establishing rigorous stability guarantees for systems with state-dependent generation of consolidation content remains an open theoretical problem (\S 8).

\subsection{Architectural Safeguards}

Four architectural properties support bounded tracking, each addressing a distinct failure mode.

\begin{enumerate}[nosep]
    \item \textbf{Sufficient dream samples per consolidation cycle.} Variance reduction in the gradient estimate. A single-timescale SA update uses a Monte Carlo estimate of the expected gradient under the current state-dependent sampling distribution; the estimation quality depends on sample size. Insufficient samples produce high-variance updates that can destabilize the iterate even under an otherwise stable mean ODE. The threshold for ``sufficient'' is empirical and will depend on deployment characteristics; the simulation in \S 4.4 establishes what the threshold looks like in the toy case.
    \item \textbf{Norm-bounded trust region.} Magnitude control independent of rank. The rank constraint limits directional complexity but not step magnitude; a rank-$r$ matrix can have arbitrarily large operator norm. The trust region caps the operator norm of each factor gradient, bounding per-cycle factor movement. As \S 2.4 states, this is not a formal compactness guarantee for the induced deformation; empirically, no safeguarded seed diverged in any simulated regime (\S 4.4), and the formal gap is flagged in \S 8.2.
    \item \textbf{Shadow adapter (Polyak--Ruppert averaging).} Readout variance reduction, and rotational damping where rotation arises. The averaged iterate concentrates much more tightly around the attractor than the live iterate, and its dispersion is nearly independent of the live stepsize at slow averaging rates (\S 4.4); the state the system reads out as its stable representation therefore enjoys a substantially tighter bounded-tracking guarantee than the live trajectory it is computed from. Separately, co-constitutive systems with state-dependent noise can in principle exhibit rotational dynamics analogous to those observed in GAN training \cite{mescheder2018}: an update that is correct on average moves the mean gradient elsewhere, producing limit cycles. Iterate averaging damps such dynamics by improving effective contractivity of the update map \cite{mou2020}. The first role is established empirically in the simulation; the second requires gradient structure with opposing components that a single descent objective cannot produce, and remains theoretically motivated rather than tested (\S 4.4).
    \item \textbf{Rank constraint as capacity bound.} Directional complexity limit. Rank limits the number of directions in which each update can move, ensuring that consolidation never closes more gradient dimensions than the surrounding topology can absorb. This prevents the failure mode of attempting to encode everything at once.
\end{enumerate}

These four safeguards are not redundant. A norm bound prevents blow-up but does not tighten the readout. Shadow averaging tightens the readout and damps rotation where it arises, but does not limit directional complexity---and, having no feedback into the live trajectory, it cannot substitute for the safeguards that do. Rank limits complexity but not magnitude. Sufficient dream samples ensure estimation quality but do not constrain the update itself. The bounded-tracking framework identifies the joint role of all four: together they produce the conditions under which iterates concentrate around a stable set rather than diverging or cycling.

\subsection{Empirical Validation}

The bounded-tracking thesis is supported by a completed program of toy simulations of the consolidation dynamics under conditions where the architecture's safeguards can be individually ablated. The program spans three adapter scales ($d = 8, r = 2$; $d = 32, r = 4$; $d = 128, r = 8$), fixed and state-dependent dream distributions (dream covariance $\Sigma = \sigma^2 I + \gamma\, \Delta W^\top \Delta W$, the minimal instantiation of the \S 5.6 coupling), the \S 5.4 symmetry-breaking sampler, and a sweep of shadow averaging rates. All runs use a persistent target perturbation, which provides the non-vanishing gradient signal at convergence that bounded tracking requires. The iterate exhibits the predicted structure in every non-pathological configuration---transient convergence followed by stable bounded fluctuation around an attractor---with live dispersion scaling as $\alpha^p$, $p \approx 0.5$ at small stepsize and dimension (matching constant-stepsize SA theory), steepening at large $\alpha$ and $d = 128$ to $p \approx 0.6$--$0.7$ in the fixed-distribution regime and as far as $p \approx 0.8$ under state-dependent dreams (with local slopes near $0.9$ at $d = 32$, $\gamma = 3$)---a quantitative deviation from the linear-theory exponent that we report without explanation (\S 8.2). The shadow readout disperses substantially less than the live iterate in every regime tested, and at slow averaging rates its dispersion is nearly independent of $\alpha$: averaging effectively decouples the dispersion of the system's stable self-representation from the live stepsize. Among the four safeguards, the trust region is load-bearing against blowup (with the clip removed, every seed at $d \geq 32$ diverges; note that this ablation is run under an added stress schedule---10$\times$-amplified inputs every 50 cycles---that the safeguarded baseline does not receive, so it demonstrates the unclipped update's vulnerability under input stress rather than a pure single-variable removal; an unconfounded rerun is flagged in \S 8.2); insufficient sample averaging inflates dispersion by 1.3--1.5$\times$ at $d = 32$ and in the fixed-distribution $d = 128$ regime, by 1.55--1.65$\times$ at $d = 8$, and by 1.9--2.75$\times$ under state-dependent dreams at $d = 128$; and the rank constraint's role is regime-dependent in an informative way---under state-dependent dreams the unconstrained adapter self-regularizes to effective rank $r$, because dream noise confined to the adapter's receptive field generates no gradient signal outside the adapter's own subspace (\S 5.6). The shadow ablation cannot be performed as originally conceived: the shadow average has no feedback into the live trajectory by construction, so removing it leaves the live trajectory identical, and the informative measure of its contribution is the ratio of live to shadow readout dispersion---which ranges from roughly $1.6\times$ at the conservative default rate in the matched regime ($1.2$--$1.5\times$ in other tiers) to roughly $9\times$ at slow rates. Detailed numerical results are reported in \S 8.

The simulation is a targeted test of the stability framework, not a substitute for empirical validation of the architecture as a whole. It establishes that bounded tracking obtains where the framework's assumptions hold, characterizes which safeguards are load-bearing in which regimes, and separates the shadow adapter's two roles: the readout variance reduction is measured and substantial, while the rotational damping is constitutively untestable in this setting---a single descent objective lacks the opposing-gradient structure that produces rotation, and direct diagnostics (trajectory autocorrelation, principal-component projections, spectral analysis) confirm its absence. The rotational-damping role therefore remains theoretically motivated rather than tested, pending a regime with genuine two-player structure. Extension to the full deployment on a trained transformer is a separate empirical question outlined in \S 8.

One limit of the bounded-tracking framing itself must be stated plainly. Bounded tracking guarantees that the adapter concentrates around an attractor; it does not guarantee that the attractor is good. Boundedness (this section) and usefulness (the functional questions of \S 8.1) are decoupled properties, and no signal internal to the consolidation loop links them: a consolidation loop can be stably bounded, useless, and register nothing wrong from the inside. This is the sharpest open problem the architecture faces, and it is why the external scaffolding of \S 5.7 and the external evaluation questions of \S 8 are not optional refinements but the other half of the stability story.


\section{Dream Content Generation}

The consolidation objective (\S 3) assumes access to an ensemble of dream variations but does not specify the generative mechanism that produces them. This section develops that mechanism. The generative process combines token-space context variation with embedding-space structured noise, and is guided by two endogenous signals derived from the adapter itself: a structural maturity signal capturing the directional breadth of adapter commitments, and an experiential maturity signal capturing cumulative consolidation history. These signals govern salience selection, noise anisotropy, and the developmental scheduling that carries the adapter from cold-start initialization through mature self-directed consolidation.

\subsection{The Adapter's Receptive Field}

A LoRA adapter for the query projection takes the form $\Delta W_Q = B_Q A_Q$, where $A_Q \in \R^{r \times d_\text{in}}$ and $B_Q \in \R^{d_\text{out} \times r}$. The effective deformation $\Delta W_Q$ is what modifies processing; the specific factorization into $B_Q$ and $A_Q$ is not unique, since multiplying $A_Q$ by any scalar $c$ and $B_Q$ by $1/c$ leaves $\Delta W_Q$ unchanged. Quantities that drive the architecture's developmental scheduling therefore depend on the effective deformation $\Delta W_Q = B_Q A_Q$ rather than on the individual factors, to ensure that the developmental clock indexes functional properties of the learned memory rather than an arbitrary factorization gauge.

Define the \emph{adapter sensitivity matrix}:
\begin{equation}
    S_Q = (B_Q A_Q)^\top (B_Q A_Q) = \Delta W_Q^\top \Delta W_Q \in \R^{d_\text{in} \times d_\text{in}},
\end{equation}
and analogously $S_V = (B_V A_V)^\top (B_V A_V)$ for the value projection. These are positive semidefinite matrices of rank at most $r$ that encode which directions in input space the effective deformation responds to and how strongly. They define the adapter's receptive field in hidden-state space and are invariant to reparameterization of the factors.

\subsection{Structural and Experiential Maturity}

Two distinct aspects of adapter development must be tracked separately. The first is \emph{structural} maturity: how much of the available rank capacity has been committed to substantive directional structure. The second is \emph{experiential} maturity: how much consolidation history has accumulated, independent of the current directional commitments. These are genuinely different axes. An adapter bootstrapped through a curriculum may be structurally mature but experientially young; a long-running adapter whose commitments have been reshaped by recent experience may be experientially mature but structurally fluid.

Define the \emph{structural maturity} signal as the participation ratio of the combined sensitivity matrix:
\begin{equation}
    \mu_\text{struct} = \frac{\left(\sum_i \sigma_i\right)^2}{r_\text{eff} \cdot \sum_i \sigma_i^2},
\end{equation}
where $\sigma_i$ are the singular values of $S_Q + S_V$, and $r_\text{eff} = \min(2r, d_\text{in})$ is the maximum effective rank of the combined sensitivity matrix. This normalization makes $\mu_\text{struct} \in [0, 1]$, with $\mu_\text{struct} = 1$ when all available directions carry equal weight. The participation ratio distinguishes an adapter that has committed strongly along few directions (specialist, low $\mu_\text{struct}$) from one that has distributed its capacity across many (broad, high $\mu_\text{struct}$), which a pure norm-based signal cannot do. The normalization by $r_\text{eff}$ rather than $d_\text{in}$ matters because $S_Q + S_V$ has rank at most $2r$; normalizing by $d_\text{in}$ would compress $\mu_\text{struct}$ into a small fraction of $[0,1]$ and prevent the scheduling mechanisms that depend on it from spanning their intended range. The Phase 1 exit threshold (\S 5.7) is calibrated against $\mu_\text{struct}$ defined this way; deployments with substantially different rank choices will require recalibration, not merely as a precaution but as a measured fact. Instrumenting both maturity signals in the \S 4.4 simulation confirms that they are monotone, bounded, and saturating as designed,\footnote{Two caveats on this instrumentation. The simulation's single-adapter analog normalizes the participation ratio by $r_\text{eff} = r$ rather than the $\min(2r, d_\text{in})$ defined here, and the noise-anisotropy \emph{scheduling} of \S 5.5 ($\sigma_\text{probe}^2/\sigma_\text{base}^2 = \gamma \cdot \mu_\text{struct}$) was not exercised anywhere in the program---the simulation couples dream covariance to the adapter at fixed $\gamma$ only, so the scheduling remains untested. We note also that monotonicity and saturation are partly built into these signals by definition (ratcheting; a saturating function of cycle count); the informative simulation output is the threshold-crossing behavior, not the monotonicity.} but the cycle at which $\mu_\text{struct}$ crosses any fixed threshold varies by orders of magnitude across adapter scales. A fixed threshold is therefore not portable across ranks; the natural replacement is a fraction-of-saturation rule (exit when $\mu_\text{struct}$ reaches a set fraction of its asymptotic value), with the caveat that estimating the asymptote online, without ground truth, is itself a nontrivial problem flagged in \S 8.

To prevent transient shrinkage of factor magnitudes from regressing $\mu_\text{struct}$, we track its high-water mark across consolidation cycles:
\begin{equation}
    \mu_\text{struct}^{(t)} = \max_{s \leq t} \mu_\text{struct}(\theta_s).
\end{equation}
This ratcheting reflects a developmental observation: structural organization, once established, does not unwind under activation fluctuation. A specialist adapter whose activation diminishes during a quiet period has not reverted to an undifferentiated state; its directional commitments persist in the factor structure, and $\mu_\text{struct}$ should reflect this. Ratcheting can in principle over-record transient peaks; whether pure high-water-mark or windowed integration is the better empirical choice is flagged as an open question in \S 8.

Define the \emph{experiential maturity} signal as normalized cumulative consolidation history:
\begin{equation}
    \mu_\text{exp} = 1 - \exp\!\left(-\frac{N_t}{N_\text{sat}}\right),
\end{equation}
where $N_t$ is the number of consolidation cycles completed and $N_\text{sat}$ is a saturation scale setting when the signal approaches 1. This gives a smooth, monotonic measure of experiential age that saturates rather than growing without bound. $N_\text{sat}$ is a hyperparameter with a principled default (the expected number of cycles for a typical deployment to reach settled operation).

The two signals govern different decisions. Structural maturity controls salience weighting (\S 5.3), noise anisotropy (\S 5.5), and the importance-weight regime (\S 5.5); these depend on what directional commitments the adapter currently has. Experiential maturity controls loss balance (\S 3.5) and phase transitions in the developmental trajectory (\S 5.7); these depend on how much experience the adapter has accumulated, which is a different question. Conflating the two into a single signal produces failure modes in both directions: decay-driven regression of a mature adapter into exploratory scheduling, and confusion between specialist and generalist adapters with similar norms.

\subsection{Salience Selection}

Not all K/V cache content is equally worth consolidating. Within-conversation prediction error (\S 3.1) provides a salience signal over token positions.

Let $\ell_t$ denote the per-token loss at position $t$ during the original conversation, and let $h_t$ denote the corresponding hidden state at the layer where the adapter is applied. Define the \emph{salience weight}:
\begin{equation}\label{eq:salience}
    w_t = \ell_t \cdot \left( (1-\mu_\text{struct}) + \mu_\text{struct} \cdot \kappa \norm{\Delta W_Q h_t}^2 \right),
\end{equation}
where $\kappa > 0$ is a scaling constant aligning the magnitudes of the two terms. When the adapter is structurally immature ($\mu_\text{struct} \approx 0$), salience reduces to prediction error alone: every surprising token is a candidate because the adapter has not yet committed to any receptive field. As structural maturity increases ($\mu_\text{struct} \to 1$), salience becomes the product of surprise and adapter engagement: content that was surprising \emph{in a region where the adapter is active} is maximally salient, representing a failure of the current deformation in its own receptive field.

This maturity dependence is essential for avoiding a cold-start pathology. At initialization $\norm{\Delta W_Q h_t}^2 = 0$ for all tokens, so a pure adapter-engagement filter would find nothing salient and produce no consolidation signal. More subtly, even after the first few updates the adapter's receptive field reflects the accident of whichever content happened to be surprising in the earliest conversations. A salience filter that immediately locks onto the adapter's receptive field would reinforce this arbitrary initial orientation through path-dependent snowballing. The structural-maturity-gated transition from broad to focused salience prevents this by ensuring that early consolidation draws from the full landscape of surprising content before the adapter's commitments narrow the selection.

Since salience weights are used only for top-$k$ ordinal selection, absolute magnitude scaling does not affect which tokens are chosen. Select the top-$k$ positions by salience weight, yielding a set $\mathcal{T} = \{t_1, \ldots, t_k\}$ (expanded to contiguous windows) and associated hidden states $\mathcal{H} = \{h_{t_1}, \ldots, h_{t_k}\}$.\footnote{The selection of contiguous subsequences versus isolated tokens is a design choice. Biological replay appears to involve temporally coherent segments \cite{wilson1994}, suggesting that $\mathcal{T}$ should comprise short contiguous windows around each salient position rather than individual token positions.}

\subsection{Context Transplantation}

The generative mechanism for dream content is \emph{context transplantation}: embedding the salient content $\mathcal{T}$ in novel surrounding contexts generated by the model itself. This produces dream sequences in which recognizable elements from recent experience appear in structurally different settings, which stress-tests whether the adapter's response to salient content generalizes across semantic surroundings.

\medskip
\noindent\textbf{Step 1: Generate candidate contexts.} Using the current model-plus-adapter state $\theta$, sample $n$ context sequences $\{c^{(1)}, \ldots, c^{(n)}\}$ from the model's generative distribution, each conditioned on a different random prefix or prompt. These contexts should be diverse; they serve as scaffolding for the transplanted salient content rather than as targets the adapter is being trained to predict (see Step 3 for how this distinction is preserved in the loss).

\medskip
\noindent\textbf{Step 2: Sampling-distribution symmetry-breaking.} Candidate contexts generated by the adapted model are shaped by the adapter's current state. A consolidation loop that trains on samples from a distribution determined by the current adapter tends toward pathological self-reinforcement: the adapter produces contexts that confirm its existing commitments, which drive gradient signals that reinforce those commitments. We apply a mild perturbation to the sampling distribution that breaks this symmetry. For each candidate $c^{(j)}$, compute the loss under both the adapted model ($L_\theta$) and the base model ($L_W$), and weight by the tempered ratio:
\begin{equation}\label{eq:sb}
    p(c^{(j)}) \propto \left(\frac{L_\theta(c^{(j)})}{L_W(c^{(j)}) + \epsilon}\right)^{1/\tau_{\text{sb}}},
\end{equation}
where $\epsilon$ is a smoothing constant and $\tau_{\text{sb}} > 1$ is a symmetry-breaking temperature (default $\tau_{\text{sb}} = 4$). The temperature compresses the signal, so that contexts where the adapted model performs worse than the base model are modestly upweighted without dominating selection. Higher $\tau_{\text{sb}}$ produces gentler correction.

The mechanism is a sampling-distribution perturbation with a tunable strength parameter, not an audit of the adapter's blind spots. The adapter cannot propose contexts it assigns near-zero generative probability to, so regions where the adapter is most resistant are never candidates regardless of reweighting strength. This invisibility of genuine blind spots to self-directed audit is not a bug to be fixed; it is a structural property of any system auditing itself through its own generative process. The architectural response is Phase 1 external curriculum (\S 5.7), which handles the class of deficiencies that endogenous mechanisms cannot reach.

The asymmetry of failure modes motivates the bias toward undercorrection. Over-aggressive correction can degrade adapter structure by consolidating adversarial, garbage, or legitimately out-of-domain content at the expense of the adapter's productive commitments. Under-aggressive correction leaves some confirmation bias in the sampling distribution, which slows progress but does not degrade structure. The gentle default of $\tau_{\text{sb}} = 4$ preserves this margin for safety. Adversarial or low-quality content is treated as out-of-scope for this mechanism---it is handled by external quality filtering at the conversation level, not by the symmetry-breaking correction.

\medskip
\noindent\textbf{Step 3: Transplant and process.} For each selected context $c^{(j)}$, construct a dream sequence by embedding the salient content $\mathcal{T}$ within $c^{(j)}$---by insertion, replacement of a segment, or interleaving. The transplanted span itself is what defines content identity across dream variations: the same span appearing in different surrounding contexts is the same content, in the operational sense relevant to the coherence loss of \S 3.4. Process the resulting chimeric sequence through the current model-plus-adapter and extract factor-space gradients $\nabla_A L_j$ and $\nabla_B L_j$ as specified in \S 3.2.

The reconstruction loss component $L_\text{recon}^{(j)}$ is computed only on tokens within and immediately adjacent to the transplanted span, not on the surrounding generated context. The surrounding context is scaffolding rather than content the adapter is being trained to predict; its only role is to vary the processing context within which the span appears. Computing reconstruction over the full chimeric sequence would punish the model for failing to predict tokens whose only coherence is local to a generated scaffold, training robustness to incoherence rather than abstraction over the salient content. The coherence loss component, which already operates only on the salient positions $\mathcal{T}$, is unaffected.

\subsection{Embedding-Space Noise and Anisotropy}

Context transplantation operates in token space (discrete, semantically meaningful). A complementary perturbation operates in embedding space (continuous, sub-symbolic). These probe different aspects of the adapter's reliability: context transplantation tests whether the adapter's response to salient content generalizes across semantic settings, while embedding noise tests whether that response is robust to representational variation at a given point in processing. Both are needed because an adapter that generalizes across contexts but is fragile to small representational shifts, or vice versa, has not been adequately consolidated.

For each dream sequence, add Gaussian noise to the hidden states with covariance structured by the adapter:
\begin{equation}\label{eq:noise}
    \tilde{h}_t = h_t + \epsilon_t, \qquad \epsilon_t \sim \mathcal{N}\!\left(0,\; \sigma_\text{base}^2 \, I + \sigma_\text{probe}^2 \, (S_Q + S_V)\right).
\end{equation}
The isotropic component ($\sigma_\text{base}^2 I$) provides diffuse exploration of directions neither adapter attends to, preserving plasticity. The structured component concentrates perturbations in the adapters' combined receptive field, stress-testing the deformation along its most sensitive directions.

\textbf{Noise anisotropy scheduling.} The ratio of structured to isotropic noise scales with structural maturity:
\begin{equation}
    \frac{\sigma_\text{probe}^2}{\sigma_\text{base}^2} = \gamma \cdot \mu_\text{struct},
\end{equation}
where $\gamma$ is a scaling constant. When $\mu_\text{struct}$ is small, noise is primarily isotropic---exploratory. As structural commitments develop, noise concentrates in the adapters' combined subspace---probing. The use of $\mu_\text{struct}$ rather than a norm-based signal ensures that noise scheduling reflects directional breadth: a narrow specialist and a broad generalist with similar Frobenius norms will have different $\mu_\text{struct}$ and correctly receive different noise regimes.

\textbf{Importance-weight regime.} Structural maturity also governs which weighting applies to candidate contexts. When $\mu_\text{struct}$ is low, the symmetry-breaking correction of Eq.~\eqref{eq:sb} applies: the sampling distribution is mildly perturbed to prevent pathological self-reinforcement while the adapter's receptive field is still forming. As $\mu_\text{struct}$ increases, the system transitions toward weighting by adapter activation magnitude:
\begin{equation}\label{eq:mature}
    p_\text{mature}(c^{(j)}) \propto \sum_{t \in c^{(j)}} \norm{\Delta W_Q \, h_t^{(j)}}^2 + \lambda \sum_{t \in c^{(j)}} \norm{\Delta W_V \, h_t^{(j)}}^2,
\end{equation}
which stress-tests existing commitments rather than perturbing the sampling distribution away from them. The transition can be implemented as a convex combination interpolated by $\mu_\text{struct}$. Because the ratcheted structural maturity does not regress under decay or transient fluctuation, the importance-weight regime does not switch between settings for spurious reasons.

\subsection{The Circularity}

The generative process is self-referential by design. The adapter's parameters define which K/V content is salient, which candidate contexts are selected, where the embedding noise concentrates, and which importance-weight regime applies. The resulting gradients update the adapter, which changes all four for the next consolidation cycle. The adapter shapes the dreams that test the adapter that update the adapter.

This circularity is the mechanism by which the system comes to encode not just what happened, but how a system with this particular history processed what happened. Under the four architectural safeguards (\S 4.3) and the bounded-tracking framework (\S 4.1), the adapter occupies an ergodic neighborhood of a self-consistent state: one that generates dreams which, when processed, produce gradient signals whose systematic component is small relative to the stepsize. This is the formal version of experience that has been absorbed into processing topology---further consolidation cycles produce only minor refinements around the current commitments, which is the healthy operating regime for a persistent learning system.

Experience that cannot be absorbed by rank-$r$ updates---because the systematic gradient signal has intrinsic dimensionality exceeding $r$---continues to generate persistent residual signal across cycles. The architectural response is the norm-bounded trust region, which ensures that any individual cycle commits a bounded amount to any direction, and the capacity constraint of rank $r$ itself, which determines how much structural complexity the adapter can absorb at all. Beyond rank-$r$, residual signal persists as bounded variability in the adapter trajectory rather than as unbounded drift.

The minimal form of this state dependence---dream covariance biased toward the adapter's receptive field---was simulated directly in the \S 4.4 program, and the coupling proves substantive rather than cosmetic, with effects that are directionally consistent across coupling strengths. Structured dream noise \emph{stabilizes} the live iterate relative to fixed-distribution dreams, because gradient signal aligned with the adapter's sensitive directions is more informative per sample; the effect strengthens monotonically with the coupling strength $\gamma$. The same rank confinement makes the rank constraint partially self-enforcing: with the constraint removed, the adapter's effective rank collapses back toward $r$ rather than growing, because noise confined to the adapter's subspace cannot generate gradient signal outside it---under state-dependent dreams, rank becomes a capacity bound rather than an actively-binding regularizer. The confinement also has costs: it caps how much the shadow average can smooth at high dimension (averaging can only operate over the directions the noise explores), and it slows convergence at small stepsizes and high dimension well beyond the $1/\alpha$ heuristic. All four effects are consequences of a single structural fact---the consolidation noise inherits the adapter's rank structure---and they are precisely the kind of effect the co-constitution argument predicts: the adapter's current state visibly shapes the statistics of its own future updates.

\subsection{Developmental Phases: The Curriculum Problem}

The dreaming mechanism described in \S 5.1--5.6 is an audit process: it takes an existing adapter, probes its receptive field, tests whether its commitments generalize, and updates them where they fail. This is the appropriate mode of operation for a mature adapter with commitments worth testing. It is not the appropriate mode of operation for a zero-initialized adapter with no commitments at all.

At initialization the LoRA adapter is zero (standard practice sets $B = 0$ so that $\Delta W = BA = 0$). The sensitivity matrices $S_Q, S_V$ are zero. Structural maturity $\mu_\text{struct}$ is zero, which is correct. Experiential maturity $\mu_\text{exp}$ is zero, also correct. The adapter's receptive field is empty. The salience weight (\S 5.3) reduces to prediction error alone; the symmetry-breaking correction (\S 5.4) returns uniform weights because the adapted and base models are identical; the structured noise component is zero. The apparatus that makes dreaming \emph{structured} rather than random collapses.

This is not merely a technical inconvenience. The first few consolidation cycles determine the adapter's initial receptive field, and all subsequent self-reinforcing dynamics build on that foundation. If the first updates are driven by the accident of whichever content happened to be surprising in the earliest conversations, the adapter develops an arbitrary initial orientation that the feedback loop then reinforces. The structural-maturity-gated salience (\S 5.3) mitigates snowballing but cannot prevent an unrepresentative initial orientation.

The structural response is a \emph{curriculum}: externally-supervised input calibrated to produce broadly representative initial organization. This is external scaffolding, not adversarial probing, and it is architecturally distinct from what the dreaming mechanism does. The division of labor follows from the limits of endogenous audit: any system auditing itself through its own generative process has principled blind spots corresponding to what its generative distribution does not reach (\S 5.4). External input handles precisely this class of deficiency. The dreaming mechanism handles refinement and stress-testing of commitments that have already been established. These are complementary architectural roles, not competing ones.

The resulting three-phase developmental trajectory for the adapter:

\textbf{Phase 1---Curriculum (initialization).} The adapter is bootstrapped through structured interactions that establish a broadly representative initial receptive field. The learning signal is external: supervised fine-tuning on a diverse corpus, a curated onboarding sequence, or the first $N$ conversations processed with standard LoRA training rather than the dreaming mechanism. The goal is not to consolidate specific experiences but to orient the adapter so that its receptive field covers a representative region of the input space.

The adapter is not starting from nothing---the base model already has a massive set of learned representations from pretraining. The curriculum phase determines which subset of that representational capacity the adapter's receptive field will initially cover. This selection is the moment at which a generic language model begins to become a particular system with a particular processing history; this is where persistent cross-session identity begins to form. We return to this framing in \S 9.

\textbf{Phase 2---Transition (early self-directed consolidation).} The dreaming mechanism begins operating with scheduling weighted toward exploration: structural maturity is modest, so salience is driven primarily by prediction error (\S 5.3), noise is primarily isotropic (\S 5.5), and the symmetry-breaking correction is active (\S 5.4). The adapter's receptive field from Phase~1 provides enough structure for the self-reinforcing dynamics to be productive rather than pathological, but the system is still mapping the landscape of what it needs to learn.

\textbf{Phase 3---Mature consolidation.} The full \S 5 apparatus operates as described: salience is adapter-weighted, noise is structured, importance sampling targets the adapter's commitments, and the system is primarily auditing and refining an established structure rather than building one from scratch.

The transition between phases is driven by the developmental signals the architecture already computes. The Phase 1 $\to$ Phase 2 transition is governed by $\mu_\text{struct}$: Phase 1 ends when $\mu_\text{struct}$ exceeds a threshold indicating that a minimally representative receptive field has been established. Because $\mu_\text{struct}$ is ratcheted, this transition is monotonic---once Phase 2 is entered, the system does not revert to Phase 1 under activation fluctuation. The Phase 2 $\to$ Phase 3 transition is governed by a combination of $\mu_\text{struct}$ (structural commitments are established) and $\mu_\text{exp}$ (enough consolidation history has accumulated to warrant auditing rather than exploring). Candidate criteria for the Phase 1 exit threshold are flagged as open questions in \S 8; a principled default would be that $\mu_\text{struct}$ reaches the level at which the top-$r$ singular values of the sensitivity matrix are above some fraction of the corresponding base-model singular values, indicating meaningful directional commitments.

The duration of Phase 1 is an open empirical question. It likely depends on the diversity of the user's early interactions and on the dimensionality of the input space the adapter needs to cover. Initialization sensitivity is essential to investigate for practical deployment; we return to it in \S 8.

\subsection{Summary: The Consolidation Algorithm}

Algorithm~\ref{alg:consolidation} summarizes the complete dreaming consolidation pipeline for Phases 2--3. During Phase 1, this algorithm does not run; the adapter is updated via standard supervised fine-tuning until $\mu_\text{struct}$ exceeds the Phase 1 exit threshold.

\begin{algorithm}[t]
\caption{Dreaming Consolidation (Phases 2--3)}
\label{alg:consolidation}
\begin{algorithmic}[1]
\REQUIRE K/V cache $\{(h_t, \ell_t)\}$, current adapter $\theta = (A_Q, B_Q, A_V, B_V)$, base model $W$, hyperparameters $k, n, \delta_{\max}, \alpha, \beta, \gamma, \kappa, \tau_\text{sb}, \lambda, \epsilon, \eta, \lambda_\text{base}, N_\text{sat}$
\STATE Compute maturity signals: $\mu_\text{struct}$ from participation ratio of $S_Q + S_V$ (ratcheted against high-water mark); $\mu_\text{exp} = 1 - \exp(-N_t / N_\text{sat})$
\STATE Salience selection (Eq.~\eqref{eq:salience} using $\mu_\text{struct}$): top-$k$ positions $\mathcal{T}$
\STATE Context generation and symmetry-breaking correction (Eq.~\eqref{eq:sb} using $\tau_\text{sb}$), interpolated with mature-regime weighting (Eq.~\eqref{eq:mature}) by $\mu_\text{struct}$: $n'$ chimeric dream sequences $d^{(j)}$
\STATE Loss weights (using $\rho$ from previous cycle, initialized $\rho_0 = 0.5$): $\lambda_\text{recon} = \mu_\text{exp} \rho / (1 + \lambda_\text{base})$, $\lambda_\text{cohere} = (1 - \mu_\text{exp} \rho + \lambda_\text{base}) / (1 + \lambda_\text{base})$
\FOR{each $d^{(j)}$}
    \STATE Add structured noise (Eq.~\eqref{eq:noise} using $\mu_\text{struct}$)
    \STATE Compute experiential losses: $L_\text{recon}^{(j)}$ (Eq.~\eqref{eq:recon}), $L_\text{cohere}^{(j)}$ (Fisher-discriminant, Eq.~\eqref{eq:cohere})
\ENDFOR
\STATE Per-dream composite loss: $L_j = \lambda_\text{recon} L_\text{recon}^{(j)} + \lambda_\text{cohere} L_\text{cohere}^{(j)}$
\STATE Extract factor-space gradients: $G_j^A = \nabla_A L_j$, $G_j^B = \nabla_B L_j$ (and analogously for $V$)
\STATE Mean factor gradients $\bar{G}^A, \bar{G}^B$; norm-clip to trust region $\delta_{\max}$
\STATE Update live adapter: $A \gets A - \alpha \bar{G}^A$, $B \gets B - \alpha \bar{G}^B$ (and analogously for $V$)
\STATE Update $\rho$ for next cycle: $\hat{\rho}$ from factor-gradient magnitudes (Eq.~\eqref{eq:rho}); $\rho \gets (1 - \alpha_\rho) \rho + \alpha_\rho \hat{\rho}$
\STATE Homeostatic decay on factors: $A \gets (1 - \eta) A$, $B \gets (1 - \eta) B$
\STATE Update shadow adapter: $\theta_\text{shadow} \gets (1 - \beta) \theta_\text{shadow} + \beta \theta$
\STATE Update structural maturity high-water mark: $\mu_\text{struct}^{(t+1)} \gets \max(\mu_\text{struct}^{(t)}, \mu_\text{struct}(\theta_{t+1}))$; increment $N_t$
\end{algorithmic}
\end{algorithm}


\section{Biological Correspondences}

The architecture's design draws on biological memory systems in places, and we describe those correspondences here. Two framing constraints govern this section. First, biology and information theory are not separate disciplines reaching different conclusions; they are different routes to the same constraints. Biological memory systems and the architecture proposed here both face the same problem---how to update parametric structure based on experience without destabilizing existing competence---and converge on similar solutions because the problem itself imposes similar requirements. Second, biomimicry narrows the probability space of designs but does not validate any particular design choice. Evolution can produce mechanisms that are good enough rather than optimal; an architectural decision motivated by biology must also be defensible on computational grounds, and we have organized the paper to ensure that the computational arguments are independently sufficient.

This section therefore retains correspondences where the biological mechanism and the computational mechanism are doing structurally analogous work, and omits correspondences where the parallel is suggestive but not load-bearing.

\textbf{Locality of modification.} The brain implements memory primarily through synaptic-strength changes localized to specific connections rather than through global rewriting of cortical processing. The LoRA adapter implements memory through low-rank deformation of attention projections rather than through retraining of base-model weights. Both impose locality as a structural feature: most of the system remains unchanged, while a small subset of parameters is the site of all experience-driven adaptation. The motivation for locality is the same in both cases. Distributed modification of the entire processing apparatus would be both expensive and destabilizing; locality allows experience to leave a structural trace without compromising the existing capabilities the trace builds on.

\textbf{Layered memory for stability and plasticity.} Both the brain and the architecture proposed here separate memory into a transient, high-resolution layer and a persistent, lower-resolution layer. In the brain, episodic experience is held briefly in hippocampal structures with high temporal precision before being slowly integrated into neocortical representations \cite{mcclelland1995}. In the architecture, the K/V cache holds the current conversation in token-level resolution, while the LoRA adapter accumulates persistent structural commitments. The functional role is the same: fast capture protects against missing recent experience, while slow integration into a more compressed representation protects against the destabilization that would follow from immediate consolidation of every experience. Within-context test-time memorization frameworks \cite{behrouz2024titans, behrouz2025atlas} share this stability-plasticity layering for the same reason; the architecture extends the layered structure across sessions.

\textbf{Homeostatic decay during consolidation.} The synaptic homeostasis hypothesis \cite{tononi2014} holds that sleep performs global synaptic downscaling that reduces absolute synaptic strengths while preserving relative differences---a renormalization that preserves the signal carried by the synaptic configuration while restoring the dynamic range available for future learning. The homeostatic decay applied to the LoRA factors after gradient updates (\S 3.4) does the same thing for the same reason. A persistent learning system that accumulates deformation without bound eventually saturates its capacity to register novel experience; uniform shrinkage of factor magnitudes preserves relative directional commitments while restoring sensitivity. This correspondence is structural rather than analogical: both systems require capacity preservation as an architectural feature for the same computational reason.

\textbf{External scaffolding for initial development.} Mammalian development includes an extended period during which the immature nervous system is unable to support self-directed learning and depends on external structuring of input. Parents do not test infants; they provide redundant, low-dimensional input calibrated to the infant's current capacities, and adjust based on observed responses. The Phase 1 curriculum (\S 5.7) plays the same architectural role for the same reason. A self-reinforcing consolidation loop has principled blind spots corresponding to what its generative distribution does not reach (\S 5.4); externally-provided structure is the only mechanism that can address this class of deficiency. The biological observation establishes that this class of deficiency is real and that external scaffolding is the response evolution converged on; the computational argument in \S 5.4 establishes why no endogenous mechanism can substitute.

\textbf{Infant sleep inversion.} One influential reading of the developmental sleep literature associates REM-dominated infant sleep with associative recombination and abstraction-related activity, and the more SWS-weighted profile of adult sleep with replay-like consolidation of specific experience \cite{tononi2014}. The self-regulating loss balance (\S 3.5) inverts the same way: a structurally mature but experientially young adapter prioritizes coherence (abstraction) over reconstruction (specific encoding), and the balance shifts toward reconstruction as experiential maturity accumulates. The biological literature here is contested in detail, and the parallel is architectural rather than one-to-one. What the parallel and the computational story share is the underlying constraint: schema formation must precede specific-experience encoding because specific experiences cannot be productively encoded into an undifferentiated representational substrate.

The correspondences omitted from this list are not oversights. Slow-wave / REM analogies for individual mechanism choices, emotional salience as a tagging signal, resistance and unconscious avoidance as motivation for sampling correction, and detailed parallels to developmental sleep architecture all suggest themselves but do not survive the test of doing independent argumentative work. The architecture is computational; the biological correspondences above are points where the structure of the computation aligns with the structure of biological mechanisms because they are solving structurally similar problems.


\section{Contributions}

This section places the contributions in relation to the Miras framework and adjacent work.

The Miras framework \cite{behrouz2025miras}, instantiated through Titans \cite{behrouz2024titans}, Atlas \cite{behrouz2025atlas}, and Memory Caching \cite{behrouz2026memcache}, has established that gradient-based test-time memorization is viable at scale: persistent memory updated via surprise-driven gradients, integrated with attention as short-term memory, achieves strong empirical performance on language modeling, long-context retrieval, and other tasks requiring extended temporal reasoning. The framework is computational and empirical, with formal capacity bounds and parallelizability guarantees that give it firm theoretical footing.

For most of its development the framework explicitly scoped itself to within-context operation; a concurrent submission extending it toward offline self-modification is discussed at the end of this section. The architecture developed here takes the framework as a starting point and extends it in three directions, with the goal of supporting cross-session memory operation.

\textbf{Cross-session persistence requires offline episodic consolidation.} Within-context consolidation in Titans, Atlas, and Memory Caching operates online: each incoming token, or each segment, produces a gradient signal that updates persistent memory weights during forward processing. Naively extending this to cross-session operation by checkpointing memory state across conversations is mechanically straightforward but produces specific architectural problems that the within-context setting does not exhibit. There is no separation between the signal that drives processing and the signal that drives consolidation; no opportunity to re-process prior experience under updated structure; and the self-reinforcing feedback between memory state and incoming gradients, bounded by context length in the within-context case, becomes unbounded across the full trajectory of accumulated experience. The architecture proposes an offline consolidation phase---structured dreaming between conversations---as the mechanism that separates processing from consolidation, provides the window for re-examining past experience under updated structure, and introduces the safeguards required to bound the feedback loop when it operates across sessions. The three-phase developmental trajectory follows from the same analysis: a persistent learning system whose sampling distribution is shaped by its own output has principled blind spots that only external input can address (\S 5.4), and in cross-session operation this is architecturally essential rather than merely useful.

\textbf{Reconstruction-family consolidation yields memory that is operative but not reusable; coherence in attentional bias supplies the missing type structure.} Our empirical program (\S 8.1) shows that consolidating experience with reconstruction-family biases---the $L_2$, $L_p$, Huber, and dot-product objectives that populate the Miras framework---produces \emph{route-bound} memory: content reliably retrievable through the pathways that consolidation trained, at ceiling on the trained-format cloze diagnostic, degrading sharply one format away and reaching chance when an independent circuit in the same model attempts to consume it. The same content supplied as context is consumable everywhere. Reconstruction, in other words, installs \emph{use}; it does not install \emph{reusability}---in the operational sense fixed in \S 1: legibility to computations beyond the trained route, demonstrated so far at the probe level, with consumption by other circuits an open prediction. The Fisher-discriminant coherence loss of \S 3.4 supplies the representational half of what is missing: it manufactures durable, probe-legible type tokens---consistent representational regions for the same content across the contexts in which it appears---an effect that replicates from linear models to a pretrained billion-parameter transformer with large effect sizes. Whether type-level structure additionally unlocks functional consumption by other circuits is, on present evidence, an open prediction rather than a result: in our testbeds it did not, and we state the conditions under which we expect it to (\S 8). The contribution is naturally formalized as an extension of the Miras signature to accommodate cross-context attentional biases, now with an empirical account of precisely which property that extension buys. Concurrent interpretability work gives the use/reusability distinction a candidate mechanistic substrate: a capacity-limited set of verbalizable representations in intermediate layers functions as a global workspace whose contents many downstream circuits can flexibly consume, while automatic processing runs outside it \cite{workspace2026}. In that vocabulary, our measurements suggest weight-consolidated content operates at the automatic tier---reliably shaping behavior without entering the workspace---whereas in-context content is workspace-native, which would explain both route-boundness and the dominance of context on token-shaped tasks; whether coherence-consolidated content reaches the workspace but fails of uptake, or never reaches it, is directly testable, and a first measurement now exists (Experiment G in the accompanying results record): on our consolidated 1.5B adapters, in-context facts broadcast strongly to answer-position mid-layer states (counterbalanced margins 0.72--1.00 across 5/5 seeds, against a copying control at or below zero in every seed) while the adapter's own contribution to the same states is indistinguishable from zero in every arm, with only a late-layer output-pathway trace---no-broadcast rather than broadcast-without-uptake, and coherence training does not change this. On present evidence, then, weight-consolidated content operates at the automatic tier outright; the coherence term manufactures probe-legible structure without buying workspace entry. This is a single instrument on a single model at one scale, and the workspace account it presupposes is concurrent work; but as a first mechanistic completion of the behavioral picture, route-boundness, format-boundness, and compositional chance all follow from the one measured fact.

\textbf{Bounded tracking is the appropriate stability framing for persistent learning systems.} Titans and Atlas demonstrate empirical stability through careful mechanism design---gating, forgetting, momentum, weight decay---without a theoretical account of what stability is meant to mean in this setting. Section 4 develops the bounded-tracking framework from constant-stepsize stochastic approximation with Markovian noise \cite{borkar2008, mou2020}, arguing that the appropriate target for any persistent learning system is not fixed-point convergence but ergodic concentration around stable attractors. The framing is offered for Titans, Atlas, Memory Caching, and the architecture developed here alike: a language for reasoning about why the gating and forgetting mechanisms these architectures rely on should be expected to work, and what specific failure modes they collectively address. For the cross-session extension, the framework identifies four architectural safeguards---rank, norm, averaging, sample size---that jointly support bounded tracking under the additional dynamical pressure of self-generated dream content.

These three contributions are connected by the problem they jointly address. The Miras framework establishes that gradient-based persistent memory is tractable within a single context. The architecture developed here asks what additional structure is required for such memory to support cross-session learning---not retention of content across long sequences, but developmental accumulation of type-level abstractions across repeated, separately-invoked interactions, with stability properties that preserve adaptation without producing drift. The answer has three parts: offline episodic consolidation with developmental scaffolding addresses the cross-session learning problem; coherence in attentional bias addresses the cross-session symbolic operation problem; and bounded tracking addresses the stability problem that becomes acute when persistent memory operates beyond context boundaries. Section 8 outlines the empirical program that tests specific predictions from each.

\textbf{Concurrent developments in the framework.} While this paper was in revision, an anonymous submission building on the same framework lineage \cite{anonymous2026sleep} proposed a sleep paradigm for language models: between active phases, the model distills the state of its fast-updating memory into gradually expanded long-term parameters (new low-rank experts added to feedforward blocks, with previously consolidated parameters frozen), then runs an explicit \emph{dreaming} stage in which self-generated synthetic data---selected by gradient-based importance scoring and validated by reinforcement from measured task-performance improvement, following the self-adapting-models line \cite{zweiger2025seal}---drives further self-modification. The within-context scoping that this paper's first contribution argued against has thus been relaxed by the framework's own line of development, and we read the convergence in vocabulary---sleep, dreaming, NREM/REM staging---as convergence on the problem. The substantive differences track this paper's three contributions precisely. First, the sleep paradigm's answer to stability is capacity growth: consolidated parameters are frozen and new capacity is added each cycle, so stability is purchased by monotone parameter expansion rather than characterized as a dynamical property of bounded-capacity updates. The bounded-tracking framework of \S 4 is the account that fixed-capacity persistent memory requires, and we would expect it to illuminate the sleep paradigm's own inner adaptation loops as much as the adapter proposed here. Second, the sleep paradigm's dreams are kept or discarded according to an external evaluator---measured downstream task performance---which is available in benchmark settings but not in open-ended deployment between conversations. The developmental trajectory, maturity scheduling, and safeguards of \S 5 address exactly the regime where no external evaluator exists, which is the regime cross-session deployment actually occupies. Third, the sleep paradigm's objectives remain reconstruction-family throughout; cross-context representational coherence does not arise as a target. More broadly, the two proposals differ in what they take consolidation to be for: knowledge incorporation evaluated by task performance there; structural individuation of a persistent, per-deployment system here.

\textbf{Continual learning with parameter-efficient adapters.} A parallel line of work addresses interference in sequentially trained adapters by construction rather than by consolidation dynamics: O-LoRA constrains successive tasks to mutually orthogonal low-rank subspaces \cite{olora2023}, with descendants protecting the frozen weights' dominant singular subspaces \cite{oplora2025} or merging sequential updates orthogonally into a single persistent adapter \cite{mbf2025}---the closest existing systems to the per-deployment persistent adapter proposed here, and replay-free alternatives to dreaming-based consolidation. Two measured facts frame this comparison. Sequential adapter training on streaming facts causes severe interference at scale (a 71\% drop in held-out QA after $\sim$1{,}000 injected facts \cite{lin2025sparse}), which our own small-scale results qualitatively reproduce at proportionally large injection loads and our safeguard ablations partially mitigate (\S 8.1); and the low-rank constraint itself trades acquisition capacity for reduced forgetting \cite{biderman2024}, a tradeoff any adapter-as-memory proposal inherits. We read the orthogonality family and the consolidation family as addressing different halves of the problem---interference geometry versus what the memory is for---and expect deployed systems to need both.

\textbf{Other related work.} Lee et al.\ \cite{lee2026sleep} study a sleep-like consolidation mechanism in which a model performs learned offline recurrent passes that convert soon-to-be-evicted context into state-space fast weights before the K/V cache is cleared, showing that deeper offline recurrence---longer sleep---improves later reasoning over the evicted context. The mechanism is trained end-to-end and operates within a single sequence, addressing reasoning under context eviction rather than cross-session persistence; it shares with the present proposal, and with the sleep-time-compute line \cite{lin2025}, the underlying claim that converting transient context into weight-based structure is a nontrivial computation deserving its own offline phase. Dream2Learn \cite{dream2learn2026} is one of several recent works using ``dreaming'' as a label for structured generative consolidation. The shared framing---structured generation outperforming verbatim replay, biological motivation for the approach---is worth acknowledging; the mechanism, modality, and purpose differ substantively, as Dream2Learn operates in image classification with a frozen external diffusion generator and prepares a classifier for future tasks, while the architecture here uses model-internal generation to consolidate past experience into persistent parametric memory in a transformer-based setting. Active Dreaming Memory \cite{adm2025} is closer in vocabulary, with a Wake/Sleep phase architecture and biological motivation, but stores consolidated knowledge as verified semantic rules in an external database rather than in parametric memory; its ``bounded forgetting'' guarantee assumes frozen base parameters and operates over retrieval probability rather than over adapter dynamics. NeuroDream \cite{neurodream2024} introduces an explicit dream phase into neural training as a continual-learning paradigm; this is closer in mechanism family but operates as a training-time paradigm rather than as a deployed-system architecture for cross-session memory. Jeong's series on persistent latent-space memory in frozen LLMs \cite{jeong2026encdec, jeong2026deconly} addresses the same application target---cross-session persistence in frozen backbones---but through write-once memory banks injected via various adapters rather than through gradient-based consolidation of accumulated experience. The complementary learning systems framework \cite{mcclelland1995} provides the broader theoretical context for the layered memory structure shared between this architecture and the Miras framework. Continual learning approaches based on regularization \cite{kirkpatrick2017} and generative replay \cite{shin2017, tadros2022} address the same underlying problem of catastrophic forgetting through different mechanisms; the architecture here is most directly compared to generative replay, with the offline phase and self-generated content being the structural commonalities, and the integration with attentional bias and the developmental trajectory being the substantive differences. Recent work on sleep-time compute \cite{lin2025} explores the use of inference-time compute for offline processing in language models, with a different architectural target but a shared intuition that test-time compute deployment can extend beyond the single forward pass.


\section{Open Questions and Empirical Program}

The architecture has been specified with sufficient precision to produce concrete empirical predictions. Section 8.1 outlines the program of experiments that test those predictions; Section 8.2 flags deferred questions whose answers depend on running the program or that lie outside the scope of a theoretical paper.

\subsection{Committed Empirical Program}

Three experiments are central. They test the three claims (\S 7) directly, and their results determine whether the architecture's structural claims survive contact with implementation.

\textbf{Stability simulation (completed).} The bounded-tracking thesis (\S 4) is supported by a completed simulation program covering three adapter scales ($d = 8/r = 2$, $d = 32/r = 4$, $d = 128/r = 8$), fixed and state-dependent dream distributions (coupling strengths $\gamma \in \{0.3, 1, 3\}$), the \S 5.4 symmetry-breaking sampler ($\tau_\text{sb} \in \{1, 4\}$), shadow averaging rates spanning $\beta \in [5 \times 10^{-4}, 0.5]$, and trajectories up to 30{,}000 cycles where convergence demanded it. The principal results:

\begin{enumerate}[nosep]
    \item \emph{Bounded tracking obtains} in every non-pathological configuration: no seed diverges with the trust region active; with the clip removed (an ablation run under an added 10$\times$-input stress schedule, \S 4.4), every seed diverges at $d \geq 32$ (10/10 at $d = 32$, 20/20 at $d = 128$), in both fixed and state-dependent regimes. At $d = 8$ the same ablation produces bounded blow-up in the fixed regime (one of ten seeds diverged under symmetry-breaking), indicating the trust region's criticality grows with dimension.
    \item \emph{Live dispersion scales as} $\alpha^p$ with $p \approx 0.5$ at small stepsize and dimension---the constant-stepsize SA prediction---steepening at large $\alpha$ and $d = 128$ to $p \approx 0.6$--$0.7$ in the fixed regime and up to $p \approx 0.8$ under state-dependent dreams (local slopes near $0.9$ at $d = 32$, $\gamma = 3$); the steepening is unexplained (\S 8.2).
    \item \emph{The shadow readout is uniformly tighter than the live iterate}, by a factor ranging from roughly $1.6\times$ at the conservative default $\beta = 0.05$ in the matched regime ($1.2$--$1.5\times$ in other tiers) to roughly $9\times$ at $\beta = 2 \times 10^{-3}$ in stationary regimes, and its dispersion is nearly flat in $\alpha$ at small $\beta$. Averaging effectively decouples readout dispersion from the live stepsize; smoothing improves monotonically as $\beta$ decreases across $[5 \times 10^{-3}, 0.5]$ (below $\beta \approx 2 \times 10^{-3}$, finite-run burn-in dominates the measurement; \S 2.5), so the matched-timescale intuition $\beta \approx \alpha$ has no support in stationary regimes.
    \item \emph{The shadow ablation is constitutively null as originally designed}: the average has no feedback into the live trajectory, so removing it cannot change live dispersion. The informative metric is the live-to-shadow dispersion ratio above. Rotational dynamics do not arise in the toy---a single descent objective lacks the opposing-gradient structure rotation requires, and direct diagnostics confirm their absence---so the rotational-damping role of averaging is untested rather than falsified.
    \item \emph{Insufficient sample averaging} inflates dispersion $1.3$--$1.5\times$ at $d = 32$ and in the fixed $d = 128$ regime, $1.55$--$1.65\times$ at $d = 8$, and $1.9$--$2.75\times$ under state-dependent dreams at $d = 128$.
    \item \emph{State-dependent coupling is substantive} (\S 5.6): it stabilizes the live iterate ($3$--$20\%$ lower dispersion at $d = 32$, from $-3\%$ at $\gamma = 0.3$ to $-20\%$ at $\gamma = 3$, monotone in $\gamma$), drives the unconstrained adapter to self-regularize back to effective rank $r$ (where the fixed-distribution regime shows rank inflation to $17$--$20$ directions), caps the shadow's smoothing ceiling at high dimension through rank confinement, and slows convergence at small $\alpha$ and high $d$ well beyond the $1/\alpha$ heuristic.
    \item \emph{The maturity signals behave as designed} ($\mu_\text{struct}$, $\mu_\text{exp}$ monotone and saturating), with Phase 1 threshold-crossing times varying by orders of magnitude across scales---the empirical basis for the fraction-of-saturation rule recommended in \S 5.2.
\end{enumerate}

Simulation code and complete numerical results accompany the paper. The toy is targeted: it establishes the stability framework's predictions where its assumptions hold. The coherence loss is not exercised---the linear toy has no notion of content identity across contexts---and rotation requires two-player gradient structure the toy cannot produce; both are noted below. Whether the results extend to the full architecture on a trained transformer is the next empirical step.

\textbf{Cross-session benchmark against an online consolidation baseline.} The central empirical question is whether the offline consolidation mechanism supports cross-session learning that an online consolidation mechanism does not. A controlled comparison would deploy an online checkpointed gradient-based memory baseline---ideally a Titans implementation if a public version becomes available, a continuously-trained LoRA adapter as a simpler proxy that captures the relevant online-vs-offline contrast, or a task-reward-validated self-adaptation baseline from the SEAL/sleep line \cite{zweiger2025seal, anonymous2026sleep}, which would additionally isolate the contribution of operating without an external evaluator---and the Dreaming LoRA architecture on a benchmark requiring retention across separate conversations: structural consistency tasks (does the system maintain the same abstraction of a concept introduced in conversation A when queried in conversation B), accumulation tasks (does the system integrate information presented across multiple conversations rather than only within each one), and cold-start tests (does the system exhibit the path-dependent pathology in initial conversations that Phase 1 curriculum is meant to address). The prediction is that offline consolidation produces measurable advantages on tasks where structural revision matters, while online consolidation may be sufficient where only retention is required. The choice of baseline architecture affects what the comparison demonstrates rigorously but not whether the comparison is meaningful: the contrast that matters is between consolidation regimes, not between specific architectures.

\textbf{Coherence loss validation (first results).} A controlled program comparing reconstruction-only against composite consolidation has now been run at three scales: a linear content/context model, a from-scratch toy transformer with an entity-attribute grammar, and a pretrained 1.5B-parameter instruct model with facts expressed across eight English registers (code and complete results accompany the paper). Four findings emerge; we note for each the testbeds that carry it. (i) \emph{Type-token formation} (all three testbeds): the composite loss produces durable cross-context identity representations---on the pretrained model, held-out-register entity retrieval 0.92 versus 0.57 and probe accuracy 0.98 versus 0.72 against reconstruction-only, with the training objective essentially unharmed (dream-domain reconstruction cost at most 6.2\% in the worst toy-grid cell; pretrained-model perplexity 17.9 $\to$ 18.1); reconstruction alone recovers roughly half the effect incidentally. A shuffled-label control shows the gains substantially require the true content structure. (ii) \emph{Route-boundness} (toy transformer and pretrained model; the linear testbed cannot express it): reconstruction-consolidated facts are retrievable at ceiling within the trained surface form (a restricted-choice cloze diagnostic over 12 invented entities, recall 1.00), degrade sharply one format away (simple QA 0.24--0.36 for city/occupation, against restricted-choice floors of 0.17/0.13), and reach chance when an independent competent circuit attempts to compose them---while the same facts supplied in context score 0.95--0.99 throughout. Consolidated memory at these scales is route-bound: operative in its trained pathways, not consumable as tokens. We read route-boundness as the signature of \emph{dispositional} memory; that reading is an interpretation, and the in-flight experiment on demonstrated-never-stated behavioral profiles is designed to test it directly on dispositional content. (iii) The coherence term does \emph{not}, at these scales, make consolidated facts consumable by other circuits: composite-consolidated facts remain compositionally inaccessible. The functional benefit of type structure is therefore an explicit open prediction, with the interaction expected where consuming circuits are rich and the consolidated content is dispositional rather than propositional. (iv) \emph{Dream diversity amplifies rather than replaces the coherence term} (linear and toy-transformer testbeds): the reconstruction-versus-composite gap grows with the number of consolidation contexts, and a directional deconfound (3 seeds, one grammar) indicates pretraining context diversity and dream context diversity act as substitutes in providing the invariant scaffold new content rides on. Separately, the safeguard ablations attached to this program measured the forgetting side: homeostatic decay is the efficient mitigation (recovering roughly a third of interference-induced loss at no cost to in-domain acquisition; its cost falls on cross-context transfer), the calibrated trust region does not move forgetting (its role is divergence prevention---the two failure modes are distinct), and interference magnitude is strongly scale-dependent in the injected-content load, consistent with measurements at larger scale \cite{lin2025sparse, biderman2024}.

Two further investigations follow naturally from the central three. They are empirical questions that the program is designed to surface answers to, even when they are not the central tests.

\textbf{Initialization sensitivity.} How quickly the adapter's receptive field becomes representative enough for self-directed consolidation to be reliable is essential for practical deployment. The Phase 1 curriculum design---how diverse it must be, how many conversations it requires, whether a shared initialization across users is preferable to per-user onboarding---is unresolved. The cross-session benchmark naturally generates data on this question by varying curriculum length and content diversity and measuring when the Phase 1 exit criterion is reached and whether subsequent self-directed consolidation produces stable structure.

\textbf{Maturity signal calibration.} The participation-ratio definition of $\mu_\text{struct}$ (\S 5.2) is one defensible choice; whether high-water-mark ratcheting is preferable to windowed integration and what saturation scale $N_\text{sat}$ to use for $\mu_\text{exp}$ remain empirical questions. The simulation has partially answered the threshold question: fixed thresholds on $\mu_\text{struct}$ are not portable across adapter scales (\S 5.2), so calibration should target a fraction of the signal's asymptotic value, and the open problem becomes online estimation of that asymptote in a deployed system without ground truth.

\subsection{Deferred Questions}

The following items are real but appropriately deferred from the present paper. They are flagged here for completeness without being developed further.

\textit{Composite loss calibration.} The Fisher-discriminant smoothing $\epsilon$, homeostatic decay rate $\eta$, symmetry-breaking temperature $\tau_\text{sb}$, and coherence floor $\lambda_\text{base}$ have principled defaults; their interactions and optimal values are empirical.

\textit{Reconstruction pressure scaling.} The current $\rho$ definition (\S 3.5) saturates near 1 for most realistic gradient norms. If the simulation reveals insufficient variation, the natural extension is to normalize against a running reference scale (low-quantile or median of recent gradient norms) rather than the constant $\epsilon$, with attention to the temporal coupling this introduces alongside the existing one-cycle lag.

\textit{Coherence loss diagnostics.} The Fisher-discriminant ratio could in principle be improved by artificial separation among salient content rather than genuine abstraction. The composite-loss balance and residual reconstruction pressure should jointly prevent this, but the failure mode warrants empirical monitoring; if diagnostics flag it, candidate mitigations include stop-gradient on the denominator or capping its effect.

\textit{Stability formalization and unconfounded ablations.} Three items from \S\S 2.4 and 4.4: a trust-region formulation that provably bounds the induced $\Delta W$ step (or an explicit factor-norm constraint) to close the compactness gap; a rerun of the trust-region ablation without the accompanying amplified-input stress schedule, to isolate the clip's contribution; and an account of the dispersion-exponent steepening beyond the linear-theory $\alpha^{1/2}$ at large stepsize and dimension.

\textit{Rotational regime.} The rotational dynamics that motivate the shadow adapter's damping role require gradient structure with opposing components, which neither the \S 4.4 toy nor any single-objective consolidation loop can produce. Testing that role requires a regime with genuine two-player structure---for instance, adapter and dream sampler optimized on decoupled timescales with partially adversarial objectives. This is a separate investigation from the present empirical program.

\textit{Multi-layer scheduling.} Different transformer layers may benefit from different consolidation regimes (reconstruction-dominant for lower layers, coherence-dominant for higher layers). Layer-specific scheduling is unexplored.

\textit{Privacy and adversarial robustness.} A per-user LoRA adapter encodes experiential residue that may be extractable through adversarial probing, and two structural points deserve statement now rather than at deployment. First, because salience is surprise-driven, an interlocutor partially controls what the system consolidates: high-surprise adversarial input is preferentially selected, which makes the salience channel itself the attack surface rather than an incidental leak; quality filtering external to the consolidation loop is therefore load-bearing, not optional. Second, the separable adapter makes \emph{deletion} easy (discard the adapter) while making \emph{selective correction} hard (no identified locus for an individual memory): the architecture is erasure-compliant but not correction-compliant, and machine unlearning within an adapter is unsolved. A full threat model---extraction, poisoning-persistence, contradiction attacks on the coherence objective---is deferred to deployment-focused work.

\textit{Base model migration.} A mechanism for adapter migration across base model versions is needed for any deployment scenario, and recent work shows the mechanism class is tractable: task adapters can be ported across base models nearly data-free using synthetic data that approximates the original training distribution \cite{translora2024}, with hypernetwork variants reported in industry \cite{portal2026}. For an \emph{experiential} adapter there is no task dataset to approximate---but the architecture already contains a candidate generator: the dreaming mechanism samples the adapter's receptive field, so the consolidation process doubles as a migration-corpus generator with no added machinery. We claim the mechanism identity, not yet the fidelity: whether dream samples preserve enough of the adapter's behavioral profile to reconstitute it on a shifted base is exactly what migration experiments must test, and our own route-boundness results counsel caution here---a generator that samples familiar routes may under-cover latent commitments, rare triggers, and behaviors the old base supplied implicitly.

\textit{Theoretical scope.} The Miras framework's capacity bounds and parallelizability guarantees assume input-driven memory updates. Offline consolidation on self-generated content operates within the framework's mechanical signature but outside the regime its current theorems cover. Extending the formal theory to cover state-dependent generation of consolidation content is an open theoretical problem.


\section{Memory and Identity}

The question which underlies this paper is: what does it mean for a transformer model to have memory? Unlike the view of memory as static recollection, we present a system of memory which is fundamentally an alteration of the self over time. The problems of memory and identity are not in fact distinct problems, but the same machinery viewed from different angles. The LoRA adapter does not record what happened; it changes how the system processes future experience in light of what happened. This is the computational analog of an old observation about biological memory: you do not remember by retrieving stored records, you remember by being a system whose dynamics have been shaped by prior experience. Memory and processing are the same machinery, viewed from different temporal perspectives.

This commitment has consequences for what cross-session operation means. The base model is, in some sense, a generic language model: a vast and capable computational substrate, but one whose response to any given query is determined by the query and its immediate context. The deployed system---base model plus persistent adapter, evolved through accumulated consolidation cycles---is no longer generic. Its receptive field has been shaped by the specific arc of experience it has processed, its symbolic abstractions have been consolidated against the specific contexts in which they have appeared, and its developmental trajectory has produced commitments that subsequent experience refines but does not replace. What the system is, at any given moment, is what the consolidated history of its processing has made it.

The type-level structure that coherence consolidation builds carries this individuation. Within-context symbolic operation in the base model has this same character: any model with the same architecture and training would track entities across passages in similar ways. Cross-session symbolic operation, built into adapter structure through repeated coherence consolidation across the specific contexts the system has encountered, is not generic. The same content---a person, a project, a recurring concern---gets a representational region in this system that reflects this system's history with that content. The representations are recognizable as referring to the same thing across sessions because the coherence loss makes them so, and the specific shape of each region is a function of the specific contexts in which the content has been encountered. The system's grasp of the content is not a record of the content; it is a structural feature of how this system processes anything that touches that content. Two systems with the same base model and different adapters would understand the same words differently, because the words activate different consolidated regions.

The empirical program (\S 8.1) gives this argument a sharper computational form than we had when it was first made. What follows in this paragraph is interpretation of those measurements---conditional on the concurrent workspace account replicating, and on behavioral tests of self-access that have not yet been run---but it is the interpretation the measurements invite. Consolidation by reconstruction alone produces memory that operates dispositionally: it reliably alters how the system behaves in the pathways experience carved, while leaving the experience only weakly recognizable as an object---roughly half the identity legibility that coherence training produces, and consumable by no circuit outside the trained route in any of our tests. Such a system has been shaped by its history and shows, on our instruments, no sign of recognizing that history \emph{in} itself: influence without acquaintance. What the coherence objective adds, measurably, is the type-token layer---each consolidated content holding a stable representational region that survives contextualization and is legible beyond the route that wrote it. If identity is consolidated history, this layer is the difference between memory that merely drives a system and memory the system can be said to \emph{have}: it is what makes the history readable by the processes that continue it, including the consolidation loop itself (a requirement \S 3.4 states as untested, and the subject of the two-cycle self-consumption experiment pre-registered in \S 8). Concurrent interpretability work sharpens the stakes of this distinction: if flexible consumption and self-report run through a workspace of verbalizable representations \cite{workspace2026}, then route-bound memory is memory with no path to report: experience that shapes the system while never, on our one instrument, entering the band through which report flows. \emph{Unconscious} is the tempting word, and we decline it only because a claim of that shape deserves more than one instrument; what the measurement licenses is narrower and stranger---influence without acquaintance, held open for the mechanistic and behavioral tests that could close it. Facts---at least the token-shaped, fixed-format facts we measured---were carried better by context than by weights in every condition tested, and we conjecture the pattern is structural rather than an artifact of scale; dispositions, on that conjecture, are parametric memory's native cargo, and the coherence layer is what keeps a dispositional self legible to itself. Memory as altered perception, with an index.

This is what the curriculum problem (\S 5.7) is really about. The first phase of an adapter's life is the moment at which a generic substrate begins to become a particular system. The external scaffolding is needed precisely because the system does not yet have any structure that would allow self-directed development to produce something representative of the user's actual world rather than an arbitrary projection of accidents from the earliest interactions. After Phase 1, the system has a starting orientation. After many cycles of dreaming and waking and dreaming again, the system has a history. The persistent adapter is the residue of that history, and the system that runs through it is constituted by what that history has made of it. This is what makes cross-session memory a different kind of question than within-context retrieval. Within context, the model has access to material; across context, the model has been shaped by material. Access is retrieval; being shaped is development.

The question is what the appropriate stability target is for a system understood this way. A persistent memory that converges to a fixed adapter state has stopped developing; the bounded-tracking framing of \S 4 is not just a mathematical convenience but a substantive claim about what healthy persistent memory looks like. The system always slightly in motion, always processing, always slightly shifted by its current experience, but never undergoing catastrophic drift, is the system that continues to develop without losing what it has become. A system that converges has finished. A system that destabilizes has lost continuity. Bounded tracking is the regime in which neither happens, and it is also the regime in which something resembling persistent identity through time remains intelligible. We have proposed an architecture engineered to target this regime; whether it occupies it, beyond the toy regime in which the signature has been observed, is what the empirical program exists to determine.

The empirical work outlined in \S 8 will determine whether the proposal is tractable in practice. The framework, if it survives that contact, makes specific claims about what transformer-based persistent memory can be made to do, and provides a principled account of why the architectural choices matter. Whether and how the resulting systems should be deployed is a separate question, with consequences that exceed what any architecture paper can resolve. We have aimed for the prior question---what such systems would be like, computationally and stability-theoretically, if they existed---and offered a framework within which both the technical and the deployment-level questions can be asked precisely.


\begin{thebibliography}{99}
\small

\bibitem{behrouz2024titans} A.\ Behrouz, P.\ Zhong, and V.\ Mirrokni. Titans: Learning to memorize at test time. arXiv:2501.00663, 2024.

\bibitem{behrouz2025atlas} A.\ Behrouz et al. Atlas: Learning to optimally memorize the context at test time. arXiv:2505.23735, 2025.

\bibitem{behrouz2025miras} A.\ Behrouz et al. Test-time memorization with the Miras framework: Attentional bias as a unifying perspective. 2025.

\bibitem{behrouz2026memcache} A.\ Behrouz et al. Memory caching: RNNs with growing memory. arXiv:2602.24281, 2026.

\bibitem{anonymous2026sleep} Anonymous. Language models need sleep: Learning to self modify and consolidate memories. Under double-blind review at ICLR 2026, 2026.

\bibitem{biderman2024} D.\ Biderman et al. LoRA learns less and forgets less. \textit{TMLR}, 2024. arXiv:2405.09673.

\bibitem{borges1942} J.\ L.\ Borges. Funes el memorioso. In \textit{Ficciones}, 1942.

\bibitem{borkar2008} V.\ S.\ Borkar. \textit{Stochastic Approximation: A Dynamical Systems Viewpoint}. Springer, 2008.

\bibitem{cha2021} H.\ Cha, J.\ Lee, and J.\ Shin. Co$^2$L: Contrastive continual learning. ICCV, 2021.

\bibitem{chen2020sdrl} T.\ Chen, T.\ Diethe, and P.\ Flach. Discriminative representation loss (DRL): A more efficient approach than gradient re-projection in continual learning. arXiv:2006.11234, 2020.

\bibitem{deperrois2022} N.\ Deperrois, M.\ A.\ Petrovici, W.\ Senn, and J.\ Jordan. Learning cortical representations through perturbed and adversarial dreaming. \textit{eLife}, 11, 2022.

\bibitem{dream2learn2026} G.\ Calcagno et al. Dream2Learn: Structured generative dreaming for continual learning. arXiv:2603.01935, 2026.

\bibitem{adm2025} B.\ T.\ Tutuncuoglu. Active Dreaming Memory: Biologically-inspired episodic consolidation for lifelong learning in autonomous agents. engrXiv preprint, December 2025.

\bibitem{neurodream2024} B.\ T.\ Tutuncuoglu. NeuroDream: A sleep-inspired memory consolidation framework for artificial neural networks. SSRN preprint, December 2024.

\bibitem{jeong2026encdec} H.\ Jeong. Trained persistent memory for frozen encoder--decoder LLMs: Six architectural methods. arXiv:2603.16413, 2026.

\bibitem{jeong2026deconly} H.\ Jeong. Trained persistent memory for frozen decoder-only LLMs. arXiv:2603.22329, 2026.

\bibitem{fountas2024} Z.\ Fountas et al. Human-like episodic memory for infinite context LLMs. arXiv:2407.09450, 2024.

\bibitem{gomezvilla2024} A.\ Gomez-Villa et al. Exemplar-free continual representation learning via learnable drift compensation. ECCV, 2024. arXiv:2407.08536.

\bibitem{hinton1995} G.\ E.\ Hinton, P.\ Dayan, B.\ J.\ Frey, and R.\ M.\ Neal. The wake-sleep algorithm for unsupervised neural networks. \textit{Science}, 268:1158, 1995.

\bibitem{hu2021} E.\ J.\ Hu et al. LoRA: Low-rank adaptation of large language models. arXiv:2106.09685, 2021.

\bibitem{kirkpatrick2017} J.\ Kirkpatrick et al. Overcoming catastrophic forgetting in neural networks. \textit{PNAS}, 114:3521, 2017.

\bibitem{lin2025sparse} J.\ Lin et al. Continual learning via sparse memory finetuning. arXiv:2510.15103, 2025.

\bibitem{lauand2023} R.\ Lauand, I.\ Kontoyiannis, and S.\ Meyn. The case for and against fixed step-size. arXiv:2309.02944, 2023.

\bibitem{lee2026sleep} S.\ Lee, S.\ McLeish, T.\ Goldstein, and G.\ Fanti. Language models need sleep. arXiv:2605.26099, 2026.

\bibitem{lin2025} K.\ Lin et al. Sleep-time compute: Beyond inference scaling at test-time. arXiv:2504.13171, 2025.

\bibitem{mbf2025} Merge before forget: A unified LoRA via orthogonal sequential merging for continual learning. arXiv:2512.23017, 2025.

\bibitem{mcclelland1995} J.\ L.\ McClelland, B.\ L.\ McNaughton, and R.\ C.\ O'Reilly. Why there are complementary learning systems in the hippocampus and neocortex. \textit{Psychological Review}, 102(3):419, 1995.

\bibitem{mescheder2018} L.\ Mescheder, A.\ Geiger, and S.\ Nowozin. Which training methods for GANs do actually converge? ICML, 2018.

\bibitem{mou2020} W.\ Mou et al. Linear stochastic approximation: How far does constant step-size and iterate averaging go? AISTATS, 2020.

\bibitem{oplora2025} OPLoRA: Orthogonal projection LoRA prevents catastrophic forgetting during parameter-efficient fine-tuning. arXiv:2510.13003, 2025.

\bibitem{portal2026} Ramp Labs. PorTAL: Portable task adapters for LLMs. Industry research article, July 2026. (Unreviewed.)

\bibitem{shin2017} H.\ Shin et al. Continual learning with deep generative replay. NeurIPS, 2017.

\bibitem{tadros2022} T.\ Tadros et al. Sleep-like unsupervised replay reduces catastrophic forgetting in artificial neural networks. \textit{Nature Communications}, 13:7742, 2022.

\bibitem{workspace2026} Anthropic Interpretability Team. Verbalizable representations form a global workspace in language models. Transformer Circuits Thread, 2026. \texttt{transformer-circuits.pub/2026/workspace}.

\bibitem{translora2024} R.\ Wang et al. \textit{Trans-LoRA}: towards data-free transferable parameter efficient finetuning. arXiv:2405.17258, 2024.

\bibitem{tononi2014} G.\ Tononi and C.\ Cirelli. Sleep and the price of plasticity: From synaptic and cellular homeostasis to memory consolidation and integration. \textit{Neuron}, 81(1):12--34, 2014.

\bibitem{olora2023} X.\ Wang et al. Orthogonal subspace learning for language model continual learning (O-LoRA). Findings of EMNLP, 2023. arXiv:2310.14152.

\bibitem{wilson1994} M.\ A.\ Wilson and B.\ L.\ McNaughton. Reactivation of hippocampal ensemble memories during sleep. \textit{Science}, 265:676--679, 1994.

\bibitem{zweiger2025seal} A.\ Zweiger, J.\ Pari, H.\ Guo, E.\ Aky\"urek, Y.\ Kim, and P.\ Agrawal. Self-adapting language models. arXiv:2506.10943, 2025.

\end{thebibliography}

\end{document}
